%!TEX TS-program = pdflatex
%!TEX encoding = utf8
\documentclass[a4paper,12pt, twoside]{book}
% \usepackage[T1]{fontenc}
\usepackage[utf8]{inputenc}
\usepackage[english, spanish, es-noshorthands]{babel}
\usepackage{amsmath}
\usepackage{amssymb}
\usepackage{enumerate} 
\usepackage{fancyvrb}              % Verbatim extendido
\usepackage{makeidx}               % Índice
\usepackage{mathrsfs}
\usepackage{pgf,tikz}
% \usepackage{algorithmic}
\usetikzlibrary{decorations.pathmorphing}
\usepackage{stackrel}
\usepackage{mathpazo} 
\usepackage[scaled=.90]{helvet}
\usepackage{cmtt}
\renewcommand{\ttdefault}{cmtt}
\usepackage{a4wide}
\usepackage{appendix}
\usepackage[pdfauthor={Miguel Angel Porras},%
  pdftitle={Sistemas de reescritura},%
  pdfstartview=FitH,%
  bookmarks=false,%
  colorlinks=true,%
  urlcolor=blue,%
  unicode=true]{hyperref}
    
\usetikzlibrary{arrows}
\usetikzlibrary{decorations.pathreplacing}

\DefineVerbatimEnvironment{code*}{Verbatim}
{fontsize=\normalsize,
 fontfamily=cmtt,
 frame=single}

\DefineVerbatimEnvironment{sesion}{Verbatim}
{commandchars=\\\{\},
 codes={\catcode`$=3\catcode`$=3},
 fontsize=\normalsize,
 fontfamily=cmtt,
 fontshape=n,
 formatcom=\color{USred},
 % xleftmargin=2em % ,
 frame=leftline
}


\DefineVerbatimEnvironment{sesion*}{Verbatim}
{%commandchars=\\\{\},
 codes={\catcode`$=3\catcode`$=3},
 fontsize=\normalsize,
 fontshape=n,
 fontfamily=cmtt,
 xleftmargin=2em,
 frame=leftline}

\DefineVerbatimEnvironment{preludio}{Verbatim}
{% commandchars=\\\{\},
 codes={\catcode`$=3\catcode`$=3},
 fontsize=\normalsize,
 fontfamily=cmtt,
 frame=single,
 framesep=2mm,
 label=\mbox{Prelude}}

\DefineVerbatimEnvironment{erroneo}{Verbatim}
{% commandchars=\\\{\},
 codes={\catcode`$=3\catcode`$=3},
 fontsize=\normalsize,
 fontfamily=cmtt,
 frame=single,
 framesep=2mm,
 label=\mbox{Código erróneo}}

\DefineVerbatimEnvironment{codigo}{Verbatim}
{% commandchars=\\\{\},
 codes={\catcode`$=3\catcode`$=3},
 fontsize=\normalsize,
 fontfamily=cmtt,
  formatcom=\color{blue},
frame=single}

\DefineVerbatimEnvironment{code}{Verbatim}
{% commandchars=\\\{\},
 codes={\catcode`$=3\catcode`$=3},
 fontsize=\normalsize,
 fontfamily=cmtt,
 frame=single}

\DefineVerbatimEnvironment{nverbatim}{Verbatim}
{commandchars=\\\{\},
 codes={\catcode`$=3\catcode`$=3},
 fontsize=\normalsize,
 fontfamily=cmtt,
 xleftmargin=2em}

\DefineVerbatimEnvironment{nverbatim*}{Verbatim}
{% commandchars=\\\{\},
 codes={\catcode`$=3\catcode`$=3},
 fontsize=\normalsize,
 fontfamily=cmtt,
 xleftmargin=2em}

\newcommand{\comentario}[1]{{\colorbox{yellow}{Comentario}: \textit{#1}}}

% =====================
% = Datos importantes =
% =====================

\title{Sistemas de reescritura desde el punto de vista de la programación funcional}
\author{Miguel Ángel Porras Naranjo}
\date{\today}
\newcommand{\tutores}[1]{\newcommand{\guardatutores}{#1}}
\tutores{Prof. José A. Alonso\\
         Prof. María José Hidalgo}

% ======================
% = Páginas de títulos =
% ======================

\makeatletter
\edef\maintitle{\@title}
\renewcommand\maketitle{%
  \begin{titlepage}
      \vspace*{1.5cm}
      \parskip=0pt
      \Huge\bfseries
      \begin{center}
          \leavevmode\includegraphics[totalheight=6cm]{sello.eps}\\[2cm]
          \@title
      \end{center}
      \vspace{1cm}
      \begin{center}
          \@author
      \end{center}
  \end{titlepage}
  
  \begin{titlepage}
  \parindent=0pt
  \begin{flushleft}
  \vspace*{1.5mm}
  \setlength\baselineskip{0pt}
  \setlength\parskip{0mm}
  \begin{center}
      \leavevmode\includegraphics[totalheight=4.5cm]{sello.eps}
  \end{center}
  \end{flushleft}
  \vspace{1cm}
  \bgroup
  \Large \bfseries
  \begin{center}
  \@title
  \end{center}
  \egroup
  \vspace*{.5cm}
  \begin{center}
  \@author
  \end{center}
  \vspace*{1.8cm}
  \begin{flushright}
  \begin{minipage}{8.45cm}
      Memoria presentada como parte de los requisitos para la obtención del título de
      Grado en Matemáticas por la Universidad de Sevilla.

      \vspace*{7.5mm}

      Tutorizada por
      % \vspace*{5mm}
  \end{minipage}\par
  \begin{tabularx}{8.45cm}[b]{@{}l}
      \guardatutores
  \end{tabularx}
   \end{flushright}
      \vspace*{\fill}
   \end{titlepage}
   \pagestyle{tfg}
   \renewcommand{\chaptermark}[1]{\markright{\thechapter.\space ##1}}
   \renewcommand{\sectionmark}[1]{\markright{\thesection.\space ##1}}
   \renewcommand{\subsectionmark}[1]{}
  }
\makeatother

% ======================================
% = Color de la Universidad de Sevilla =
% ======================================

\usepackage{tikz}
\definecolor{USred}{cmyk}{0,1.00,0.65,0.34}

% =========
% = Otros =
% =========

\usepackage[]{tabularx}
\usepackage[]{enumitem}
\setlist{noitemsep}

% ==========================
% = Matemáticas y teoremas =
% ==========================

\usepackage[]{amsmath}
\usepackage[]{amsthm}
\usepackage[]{mathtools}
\usepackage[]{bm}
\usepackage[]{thmtools}
\usepackage[]{thmtools}
\newcommand{\marcador}{\vrule height 10pt depth 2pt width 2pt \hskip .5em\relax}
\newcommand{\cabeceraespecial}{%
    \color{USred}%
    \normalfont\bfseries}
\declaretheoremstyle[
    spaceabove=\medskipamount,
    spacebelow=\medskipamount,
    headfont=\cabeceraespecial\marcador,
    notefont=\cabeceraespecial,
    notebraces={(}{)},
    bodyfont=\normalfont\itshape,
    postheadspace=1em,
    numberwithin=chapter,
    headindent=0pt,
    headpunct={.}
    ]{importante}
\declaretheoremstyle[
    spaceabove=\medskipamount,
    spacebelow=\medskipamount,
    headfont=\normalfont\itshape\color{USred},
    notefont=\normalfont,
    notebraces={(}{)},
    bodyfont=\normalfont,
    postheadspace=1em,
    numberwithin=chapter,
    headindent=0pt,
    headpunct={.}
    ]{normal}
\declaretheoremstyle[
    spaceabove=\medskipamount,
    spacebelow=\medskipamount,
    headfont=\normalfont\itshape\color{USred},
    notefont=\normalfont,
    notebraces={(}{)},
    bodyfont=\normalfont,
    postheadspace=1em,
    headindent=0pt,
    headpunct={.},
    numbered=no,
    qed=\color{USred}\marcador
    ]{demostracion}
\declaretheoremstyle[
spaceabove=\medskipamount,
spacebelow=\medskipamount,
headfont=\cabeceraespecial,
notefont=\cabeceraespecial,
notebraces={(}{)},
bodyfont=\normalfont,
postheadspace=1em,
numberwithin=chapter,
headindent=0pt,
headpunct={.}
]{nota}


\declaretheorem[name=Observaci\'on, style=normal]{obse}
\declaretheorem[name=Corolario, style=normal]{coro}
\declaretheorem[name=Notaci\'on, style=nota]{nota}
\declaretheorem[name=Ejemplo, style=nota]{ejem}
\declaretheorem[name=Proposici\'on, style=normal]{prop}
\declaretheorem[name=Lema, style=normal]{lema}

\declaretheorem[name=Teorema, style=importante]{teor}
\declaretheorem[name=Definici\'on, style=importante]{defi}

\let\proof=\undefined
\declaretheorem[name=Demostraci\'on, style=demostracion]{demo}


% ============================
% = Composición de la página =
% ============================

\usepackage[
        a4paper,
    textwidth=80ex,
]{geometry}

\linespread{1.069}
\parskip=10pt plus 1pt minus .5pt
\frenchspacing
% \raggedright

% ==============================
% = Composición de los títulos =
% ==============================

\usepackage[explicit]{titlesec}

\newcommand{\hsp}{\hspace{20pt}}
%\titleformat{\chapter}[hang]
%    {\Huge\sffamily\bfseries}
%    {\thechapter\hsp\textcolor{USred}{\vrule width 2pt}\hsp}{0pt}
%    {#1}
%\titleformat{\section}
%  {\normalfont\Large\sffamily\bfseries}{\thesection\space\space}
%  {1ex}
%  {#1}
%\titleformat{\subsection}
%  {\normalfont\large\sffamily}{\thesubsection\space\space}
%  {1ex}
%  {#1}
  
% Abreviaturas matemáticas:
% -------------------------
\newcommand{\Pos}{\mbox{Pos }}
\newcommand{\Var}{\mbox{Var }}
\newcommand{\Ran}{\mbox{Ran }}
\newcommand{\VRan}{\mbox{VRan }}
\newcommand{\Sub}{\mbox{Sub }}
\newcommand{\Dom}{\mbox{Dom }}
\newcommand{\Unify}{\mbox{\texttt{Unificar}}}
\providecommand{\abs}[1]{\lvert#1\rvert}

% =======================
% = Cabeceras de página =
% =======================

\usepackage[]{fancyhdr}
\usepackage[]{emptypage}
\fancypagestyle{plain}{%
    \fancyhf{}%
    \renewcommand{\headrulewidth}{0pt}
    \renewcommand{\footrulewidth}{0pt}
}
\fancypagestyle{tfg}{%
    \fancyhf{}%
    \renewcommand{\headrulewidth}{0pt}
    \renewcommand{\footrulewidth}{0pt}
    \fancyhead[LE]{{\normalsize\color{USred}\bfseries\thepage}\quad
                    \small\textsc{\MakeLowercase{\maintitle}}}
    \fancyhead[RO]{\small\textsc{\MakeLowercase{\rightmark}}%
                    \quad{\normalsize\bfseries\color{USred}\thepage}}%
                    }

\makeatletter
\newcommand{\dotminus}{\mathbin{\text{\@dotminus}}}

\newcommand{\@dotminus}{%
  \ooalign{\hidewidth\raise1ex\hbox{.}\hidewidth\cr$\m@th-$\cr}%
}
\makeatother

%%%%%%%%%%%%%%%%%%%%%%%%%%%%%%%%%%%%%%%%%%%%%%%%%%%%%%%%%%%%%%%%%%%%%%%%%%%%%%%
%%  Construcción del índice                                                 %%
%%%%%%%%%%%%%%%%%%%%%%%%%%%%%%%%%%%%%%%%%%%%%%%%%%%%%%%%%%%%%%%%%%%%%%%%%%%%%%%

\makeindex

% =============================
% = El documento empieza aquí =
% =============================

% \includeonly{Formalizacion_del_concepto_de_terminos}

\begin{document}
\maketitle

\frontmatter
\tableofcontents

\mainmatter

\chapter*{Introducción}
\addcontentsline{toc}{section}{Introducción}

La \textbf{reescritura} es una técnica que consiste en reemplazar términos de
una expresión con \textbf{términos equivalentes}. Por ejemplo, dadas las reglas
$x * 0 \Rightarrow 0$ y $x + 0 \Rightarrow 0$ pueden usarse para
simplificar la siguiente expresión:

\[
x + (\underline{x*0}) \longrightarrow \underline{x + 0} \longrightarrow x
\]

Pero no sólo se puede aplicar para ecuaciones. Se puede usar para
demostrar enunciados. Por ejemplo, dadas las reglas
$s(M) \leq s(N) \Rightarrow M \leq N$ y
$0 \leq N \Rightarrow \texttt{True}$ pueden aplicarse a:

\[
  s(0) \leq s(x) \longrightarrow 0 \leq s(x) \longrightarrow
  \texttt{True}
\]

Donde $s(x)$ es la función sucesor de $x$.

En definitiva, podemos aplicar reescritura a cualquier tipo de
expresión (ya sea lógica, ecuacional, etc) siempre y cuando se pueda
hacer una \textbf{sustitución}. Esta sustitución se hace de término a término.

Sin embargo, debemos formalizar qué entendemos por término. Un \textbf{término}
estará construido por \textbf{variables, símbolos constantes y símbolos de
funciones}. Al estar constituidos por funciones, una representación
natural de los términos se realiza mediante árboles. Por ejemplo el
término $f(i(e),f(x,e))$ se representa como,

\begin{figure}[h]
  \centering
  \begin{tikzpicture}
    \node [circle,draw]{$f$} child {node
      [circle,draw]{$i$}
        child{node[circle,draw]{$e$}}} child {node
      [circle,draw]{$f$} child {node
        [circle,draw]{$x$}} child {node
        [circle,draw]{$e$}} };
  \end{tikzpicture}
\end{figure}

Donde $i,f$ representa un símbolo de función de aridad 1 y 2
respectivamente, $e$ representa un símbolo constante y $x$ representa
una variable.

Los anteriores ejemplos también se pueden escribir como términos,

\[
  \begin{array}{rcl}
    x +(x*0) & \longrightarrow  & +(x,*(x,0)) \\ \\
    s(0) \leq s(x) & \longrightarrow & \leq(s(0),s(x)) \\ \\

  \end{array} 
\]

Un \textbf{sistema de reescritura de términos} es un conjunto de sustituciones
término a término. A estas sustituciones las llamaremos reglas.

Estos conjuntos de reglas determinan ciertas propiedades según los
elementos que contengan. Por ejemplo,

\begin{figure}[h]
  \centering
  \begin{tikzpicture}[sibling distance=10em]
    \node [draw]{$a^{2*0} + b*0$}
    child {
      node[draw]{$a^0 + b*0$}
      child{node[draw]{$1 + b*0$}}
      child{node[draw]{$a^0 + 0$}}
      }
      child {
        node [draw]{$a^{2*0} + 0$}
        child {node [draw]{$a^0 + 0$}}
        child {node [draw]{$a^{2*0}$}} };
  \end{tikzpicture}
\end{figure}

Este diagrama de árbol representa para cada rama una aplicación de las
distintas reglas. Podemos hacernos dos preguntas fundamentales, si
este árbol será \textbf{finito} y si \textbf{todas las hojas son la misma}; es decir,
siguiendo cualquier rama, llegar hasta el mismo término. Estas dos
preguntas son los principales problemas de la reescritura.

Analizar si este árbol es finito es equivalente a preguntarse si en
algún momento podemos dejar de aplicar las reglas. Por ejemplo, para
el término $s(1) + s(0)$ y para la regla $s(3) \Rightarrow 4$, no
podemos seguir sustituyendo términos, por tanto diremos que el sistema
\textbf{termina}.

Es importante observar que no siempre tener un mayor número de reglas
es mejor. Por ejemplo, para las reglas $ x + 0 \Rightarrow x$ y
$x \Rightarrow x + 0$;

\[
  x + 0 \longrightarrow x \longrightarrow x+0 \longrightarrow \dots
\]

Otro de los problemas que tenemos que tratar es el de ver si el árbol
termina en una misma hoja. Si esta situación ocurre en nuestro sistema
diremos que es \textbf{confluente}.

Para analizar si un sistema es confluente o no debemos estudiar el
momento en el que se crea una ramificación del árbol. La idea
intuitiva es ver si la sustitución que se hace en cada rama se solapa
entre sí. A estos puntos los llamaremos \textbf{pares críticos}.

También podemos modificar el conjunto reglas de los reescritura para
que se verifiquen estas dos propiedades. A este proceso se le llama
\textbf{completación}. Los cambios que se realizan al conjunto de reglas van
desde eliminar alguna regla, hasta añadir propias.

Como es de esperar, el algoritmo de completación debe tener en cuenta
los pares críticos. Una primera idea de implementación sería modificar
las reglas hasta no tener ningún par crítico.

Sin embargo el cálculo de pares críticos es muy costoso si hablamos de
grandes sistemas de reescritura. Por ello introduciremos \textbf{el algoritmo
de Huet}, que pese a usar pares críticos en el proceso, reduce la
cantidad de cálculos a realizar.

A continuación, describiré brevemente los contenidos de cada capítulo.

En el capítulo 1, introduciremos los \textbf{sistemas de reducción abstractos} y
definiremos algunas de sus características fundamentales, algunas de
ellas se volverán a tratar con mayor profundidad en capítulos
posteriores.

En el capítulo 2, estudiaremos las definiciones de \textbf{términos y
sustituciones}, así como sus características. Además, relacionando
términos y sustituciones, daremos una caracterización de
$\xleftrightarrow{*}$,

En el capítulo 3, formalizaremos la definición de \textbf{sistema de
reescritura de términos} y estudiaremos el problema de \textbf{unificación} para
términos.

En el capítulo 4, resolveremos el problema de decidir si un sistema de
reescritura es \textbf{terminante}. Veremos que, en general, el problema es
indecidible. Por ello, le daremos varios enfoques según las
propiedades que tenga el sistema para poder decidir la terminación.

En el capítulo 5, estudiaremos la \textbf{confluencia} para los sistemas de
reescritura. Probaremos que el problema de decidir si un sistema es
confluente es indecidible. Sin embargo, a partir del estudio de pares
críticos y ortogonalidad, podremos dar varios resultados sobre la
confluencia.

Por último, en el capítulo 6, enunciaremos dos algoritmos de
\textbf{completación}. Para ambos, estudiaremos su funcionamiento y si son
correctos.

%%% Local Variables:
%%% mode: latex
%%% TeX-master: "SRT_en_Haskell"
%%% End:

\chapter{Sistemas de reducción abstractos}

En este capítulo, definiremos que entendemos como un sistema de
reducción abstracto, y a partir de este, daremos la definición de
orden, que sera imprescindible para capítulos posteriores. También
veremos en qué consiste la inducción bien fundada, el orden
lexicográfico y un poco de teoría de multiconjuntos. Estas dos últimas
incluirán una implementación en Haskell que se puede encontrar en
\texttt{Orden.hs}


\section{Equivalencia y reducción}

El objetivo principal de esta sección es definir toda la estructura de
los sistemas de reducción abstractos.

\begin{defi} 
  Un \textbf{sistema de reducción abstracto} es un par $(A, \rightarrow)$ donde
  la reducción $\rightarrow$, es una relación en el conjunto $A$; es decir,
  $\rightarrow \subseteq A \times A$.
\end{defi}

Para simplificar la notación, si $(a,b) \in \rightarrow$, escribiremos
$a\rightarrow b$.  Un sistema de reducción abstracto se puede entender
como un grafo dirigido. Por ejemplo,

\begin{figure}[h]
          \centering
          \begin{tikzpicture}
          \begin{scope}[execute at begin node=$, execute at end node=$]
          \node (v1) at (0,0) {a};
          \node (v2) at (3,0) {b};
          \node (v3) at (6,1) {c};
          \node (v4) at (6,-1) {d};
          \end{scope}
          
          \draw [-latex](v2) (v1) edge  node [auto] {} (v2);
          \draw [-latex](v3) (v2) edge  node [auto] {} (v3);
          \draw [-latex](v4) (v2) edge  node [auto] {} (v4);
          \draw [-latex](v1) (v2) edge [bend right] node [auto] {} (v1);
          \end{tikzpicture}
\end{figure}

\begin{defi} 
  Diremos que $b$ es una \textbf{reducción} de $a$ si
  $a\rightarrow a_0 \rightarrow a_1\rightarrow \dots\rightarrow b$.
\end{defi}

En el ejemplo anterior, $c$ es una reducción de $a$. Pero $d$ no es
una reducción de $c$

\begin{defi} 
  Diremos que $a$ y $b$ son \textbf{equivalentes}, y lo denotaremos por
  $a \xleftrightarrow{*} b$, si $b$ es una reducción de $a$ y $a$ es una
  reducción de $b$.
\end{defi}

En el ejemplo, $a$ y $b$ son equivalentes.
También trabajaremos con el concepto de composición.

\begin{defi} 
  Dadas dos relaciones $R \subseteq A\times B$, y $S \times B \times C$,
  definimos su \textbf{composición} como
  \[R \circ S := 
    \left\lbrace (x,z)\in A \times C | \exists y \in B. (x,y) \in R \wedge (y,z) \in S \right\rbrace   \]
\end{defi}

A partir de una reducción $\rightarrow$ se pueden calcular diferentes clausuras,

\begin{defi} 
  A partir de la definición de reducción, consideraremos las siguientes nociones, \\
  $\begin{array}{llll}
     \xrightarrow{0}         & := & \left\lbrace (x,x) | x \in A \right\rbrace   & $\textbf{Identidad}$  \\ 
     \xrightarrow{i+1}       & := & \xrightarrow{i}$ ó $ \rightarrow & \textbf{(i+1)}$\textbf{-plegado,} $  \pmb{i\geq 0} \\ 
     \xrightarrow{+}         & := & \cup_{i>0} \xrightarrow{i} & $\textbf{Clausura transitiva}$ \\ 
     \xrightarrow{*}         & := & \xrightarrow{+} \cup \xrightarrow{0} & $\textbf{Clausura transitiva reflexiva}$ \\ 
     \xrightarrow{=}         & := & \rightarrow \cup \xrightarrow{0} & $\textbf{Clausura reflexiva}$ \\ 
     \xrightarrow{-1} & := & \left\lbrace (y,x) | x\rightarrow y \right\rbrace  & $\textbf{Inversa}$ \\ 
     \leftarrow      & := & \xrightarrow{-1} & $\textbf{Inversa}$ \\ 
     \leftrightarrow         & := & \rightarrow \bigcup \leftarrow & $\textbf{Clausura simétrica}$ \\ 
     \xleftrightarrow{+}     & := & \left( \leftrightarrow \right)^{+} & $\textbf{Clausura transitiva simétrica}$ \\ 
     \xleftrightarrow{*}     & := & \left( \leftrightarrow \right)^{*} & $\textbf{Clausura reflexiva transitiva simétrica}$
   \end{array} 
   $
\end{defi}

\begin{ejem}
  A partir de la reducción:
  \begin{figure}[h]
          \centering
          \begin{tikzpicture}
          \begin{scope}[execute at begin node=$, execute at end node=$]
          \node (v1) at (0,0) {a};
          \node (v2) at (3,0) {b};
          \node (v3) at (6,1) {c};
          \node (v4) at (6,-1) {d};
          \node (v5) at (9,-2) {e};
          \end{scope}
          
          \draw [-latex](v2) (v1) edge  node [auto] {} (v2);
          \draw [-latex](v3) (v2) edge  node [auto] {} (v3);
          \draw [-latex](v4) (v2) edge  node [auto] {} (v4);
          \draw [-latex](v5) (v4) edge  node [auto] {} (v5);
          
  
          \end{tikzpicture}
  \end{figure}    
  Las relaciones según la definición anterior son,
  \begin{figure}[h]
          \centering
          \begin{tikzpicture}
          \begin{scope}[execute at begin node=$, execute at end node=$]
          \node (v1) at (0,0) {a};
          \node (v2) at (3,0) {b};
          \node (v3) at (6,1) {c};
          \node (v4) at (6,-1) {d};
          \node (v5) at (9,-2) {e};
          \end{scope}
          
          \draw [-latex](v2) (v1) edge  node [auto] {} (v2);
          \draw [-latex](v3) (v2) edge  node [auto] {} (v3);
          \draw [-latex](v4) (v2) edge  node [auto] {} (v4);
          \draw [-latex](v5) (v4) edge  node [auto] {} (v5);
          \draw [-latex](v4) (v1) edge [bend right] node [auto] {} (v4);
          \draw [-latex](v5) (v1) edge [bend right] node [auto] {} (v5);
          \draw [-latex](v3) (v1) edge [bend left] node [auto] {} (v3);
          \draw [-latex](v5) (v2) edge [bend left] node [auto] {} (v5);
          
          
          \end{tikzpicture}
          \caption{$\xrightarrow{+}$}
  \end{figure}            
  \begin{figure}[h]
          \centering
          \begin{tikzpicture}
          \begin{scope}[execute at begin node=$, execute at end node=$]
          \node (v1) at (0,0) {a};
          \node (v2) at (3,0) {b};
          \node (v3) at (6,1) {c};
          \node (v4) at (6,-1) {d};
          \node (v5) at (9,-2) {e};
          \end{scope}
          
          \draw [-latex](v2) (v1) edge  node [auto] {} (v2);
          \draw [-latex](v3) (v2) edge  node [auto] {} (v3);
          \draw [-latex](v4) (v2) edge  node [auto] {} (v4);
          \draw [-latex](v5) (v4) edge  node [auto] {} (v5);
          \draw [-latex](v4) (v1) edge [bend right] node [auto] {} (v4);
          \draw [-latex](v5) (v1) edge [bend right] node [auto] {} (v5);
          \draw [-latex](v3) (v1) edge [bend left] node [auto] {} (v3);
          \draw [-latex](v5) (v2) edge [bend left] node [auto] {} (v5);
          \draw [-latex](v2) (v1) edge [loop left] node [auto] {} (v2);
          \draw [-latex](v3) (v2) edge [loop] node [auto] {} (v3);
          \draw [-latex](v5) (v4) edge [loop below] node [auto] {} (v5);
          \draw [-latex](v5) (v3) edge [loop] node [auto] {} (v5);
          \draw [-latex](v5) (v5) edge [loop right] node [auto] {} (v5);
          
          \end{tikzpicture}
          \caption{$\xrightarrow{*}$}
  \end{figure}            
\end{ejem}

A continuación, veremos una serie de definiciones que serán
imprescindibles a la hora de estudiar los sistemas de reescritura.

\begin{defi} 
  Diremos que $P$ es la \textbf{clausura} de $R$ si $P$ es el menor conjunto
  con la propiedad $P$ que contiene a $R$.
\end{defi}

\begin{defi} 
  Diremos que $x$ es \textbf{reducible} si existe un $y$ tal que
  $x \rightarrow y$.
\end{defi}

\begin{defi} 
  Diremos que $x$ esta en su \textbf{forma normal} si $x$ no es reducible.
\end{defi}

%\begin{defi} 
%  Diremos que $\pmb{y}$ \textbf{es la forma normal de} $\pmb{x}$ si
%  $x \xrightarrow{*} y$, donde $y$ es una forma normal. Si $x$ tiene una única
%  forma normal, lo denotaremos por $x \downarrow$.
%\end{defi}

%\begin{defi} 
%  Diremos que $y$ es un \textbf{sucesor directo} de $x$ si $x \rightarrow y$.
%\end{defi}

\begin{defi} 
  Diremos que $y$ es un \textbf{sucesor} de $x$ si $x \xrightarrow{+} y$
\end{defi}      

\begin{defi} 
  Diremos que $x$ e $y$ son \textbf{confluentes} si hay algún $z$ tal que
  $x\xrightarrow{*}z \xleftarrow{*} y$. Lo denotaremos por $x \downarrow y$.
\end{defi}      

%Ejemplos 2.1.2 1)

\begin{defi} 
  Llamaremos a una reducción $\rightarrow$, \\
  $       
  \begin{array}{ll}
    $\textbf{Church-Rosser}$        & $si $ x \xleftrightarrow{*} y \Leftrightarrow x \downarrow y \\ 
    $\textbf{confluente}$   & $si $ y_1 \xleftarrow{*} x \xrightarrow{*} y_2 \Rightarrow y_1 \downarrow y_2\\ 
    $\textbf{semi-confluente}$      & $si $ y_1 \leftarrow x \xrightarrow{*} y_2 \Rightarrow y_1 \downarrow y_2\\ 
    $\textbf{terminante}$   & $si no hay una cadena infinita $ a_0 \rightarrow a_1 \rightarrow \dots \\ 
    $\textbf{normalizado}$  & $si cada elemento tiene forma normal$\\ 
    $\textbf{convergente}$  & $si es confluente y terminante$
  \end{array}     
  $
\end{defi}

La línea discontinua en los siguientes diagramas representa la
existencia de una reducción. Por ejemplo, en Church Rosser nos indica que si $ x \xleftrightarrow{*} y$, entonces $\exists z$ tal que $x \xleftrightarrow{*} z $, e $y \xleftrightarrow{*} z$

%Church Rosser
\begin{figure}[h]
  \centering
  \begin{tikzpicture}
          \node (v1) at (0,3) {x};
          \node (v2) at (3,3) {y};
          \node (v3) at (1.5,0.8) {z};
          
          \draw [latex-latex](v2) (v1) edge  node [auto] {*} (v2);
          \draw [latex-latex, dash pattern=on 2pt off 3pt on 4pt off 4pt](v3)  edge  node [auto] {*} (v1);
          \draw [latex-latex, dash pattern=on 2pt off 3pt on 4pt off 4pt](v2)  edge  node[auto] {*} (v3);
  \end{tikzpicture}
  \caption{Church Rosser}
\end{figure}

%Confluente
\begin{figure}[h]
  \centering
  \begin{tikzpicture}
    \begin{scope}[execute at begin node=$, execute at end node=$]
      \node (v3) at (0,0) {y_2};
      \node (v1) at (0,3) {x};
      \node (v4) at (3,0) {z};
      \node (v2) at (3,3) {y_1};
      \end{scope}
      
      \draw [-latex](v2) (v1) edge  node [auto] {*} (v2);
      \draw [-latex](v3) (v1) edge  node [left] {*} (v3);
      \draw [-latex, dash pattern=on 2pt off 3pt on 4pt off 4pt](v3)  edge  node [below] {*} (v4);
      \draw [-latex, dash pattern=on 2pt off 3pt on 4pt off 4pt](v2)  edge  node[auto] {*} (v4);
  \end{tikzpicture}
  \caption{Confluente}
\end{figure}

%Semi-confluente
\begin{figure}[h]
  \centering
  \begin{tikzpicture}
    \begin{scope}[execute at begin node=$, execute at end node=$]
      \node (v3) at (0,0) {y_2};
      \node (v1) at (0,3) {x};
      \node (v4) at (3,0) {z};
      \node (v2) at (3,3) {y_1};
      \end{scope}
      
      \draw [-latex](v2) (v1) edge  node [auto] {} (v2);
      \draw [-latex](v3) (v1) edge  node [left] {*} (v3);
      \draw [-latex, dash pattern=on 2pt off 3pt on 4pt off 4pt](v3)  edge  node [below] {*} (v4);
      \draw [-latex, dash pattern=on 2pt off 3pt on 4pt off 4pt](v2)  edge  node[auto] {*} (v4);
  \end{tikzpicture}
  \caption{Semi-confluente}
\end{figure}

\section{Resultados básicos}

A partir de las anteriores definiciones, vamos a demostrar una serie
de propiedades que las relacionan entre sí.

\begin{teor} \label{teor:conflu}
  Las siguientes condiciones son equivalentes,
  \begin{enumerate}
  \item $\rightarrow$ tiene la propiedad de Church--Rosser.
  \item $\rightarrow$ es confluente.
  \item $\rightarrow$ es semi--confluente.
  \end{enumerate}
\end{teor}

\begin{demo}
  La implicación $(1 \Rightarrow 2)$ se prueba suponiendo
  $y_1 \xleftarrow{*} x \xrightarrow{*} y_2$ por tanto,
  $y_1 \xleftrightarrow{*} y_2$, y por la propiedad de Church--Rosser
  $y_1 \downarrow y_2$.  

  La implicación $(2 \Rightarrow 3)$ es trivial.
%FALTA
\end{demo}

% Consecuencias directas del teorema

\begin{coro} 
  Si $\rightarrow$ es confluente y $x \xleftrightarrow{*} y$ entonces,
  \begin{enumerate}
  \item $ x \xrightarrow{*} y$, si $y$ es una forma normal. 
  \item $x=y$, si $x$ e $y$ están en su forma normal.
  \end{enumerate}         
\end{coro}

\begin{teor}
  Si $\rightarrow$ es normalizado y confluente entonces
  $ x \xleftrightarrow{*} y$ syss $\ x \downarrow = y \downarrow$. Es decir,
  cada elemento tiene una única forma normal.
\end{teor}

\section{Inducción bien fundada}

El principio de la inducción bien fundada (IBF) se trata de una generalización
de la inducción en $(\mathbb{N},>)$ para $(A,\rightarrow)$.  Formalmente, la
expresaremos como la siguiente regla:
\begin{equation*}
  \frac{\forall x \in A. (\forall y \in A. x \xrightarrow{+} y \Rightarrow P(y)) \Rightarrow P(x)} 
       {\forall x \in A. P(x)}
\end{equation*}

% INCLUIR UN EJEMPLO

\begin{teor}  
  Si $\rightarrow$ es terminante, entonces (IBF) se verifica.
\end{teor}

\begin{teor}  
  Si (IBF) se verifica, entonces $\rightarrow$ es terminante.
\end{teor}

Hemos llegado a que la definición de una relación terminante, es equivalente a
comprobar que verifica (IBF). Esta propiedad será importante a la hora de
resolver el problema de la terminación, en el capítulo \ref{cap:term}.

\begin{defi} 
  Llamaremos a una relación $\rightarrow$, \\
  $
  \begin{array}{ll}
    $
    \textbf{Ramificación finita}$   & $Si cada elemento tiene algunos sucesores finitos directos.$ \\ $
    \textbf{Globalmente finita}$    & $Si cada elemento tiene algunos sucesores finitos.$   \\ 
    $ \textbf{Acíclico}     $ & $Si no hay ningún elemento $a$, tal que $a \xrightarrow{+} a.
  \end{array}
  $
\end{defi}      

% Comentario: Precisar la definición anterior.

Vamos a probar dos resultados de esta definición. Una relación
globalmente finita es equivalente a decir que la relación
$\xrightarrow{+}$ es una ramificación finita.

\begin{lema} \label{ramfin} 
  Una relación de ramificación finita es globalmente finita si es terminante.
\end{lema}

\begin{demo}
  Usaremos la propiedad (IBF). Sea $\rightarrow$ una relación de
  ramificación finita y terminante. Vamos a probar que cada elemento
  del conjunto de los sucesores es finito. Como es cierto para cada
  sucesor directo por hipótesis, aplicando (IBF) llegamos a que su
  conjunto de sucesores es finito, y por tanto globalmente finito.
\end{demo}

\begin{lema} 
  Una relación acíclica es terminante si es globalmente finita.
\end{lema}

\section{Resultados sobre la terminación}

Uno de los métodos que estudiaremos para probar que $\rightarrow$ es
terminante, se basa en (IBF). A partir de $(A,\rightarrow)$, definiremos un
sistema de reducción abstracto terminante $(B,>)$. Definiremos una función
$\varphi : A \rightarrow B$ que debe tener la propiedad de monotonía, (ie
$x \rightarrow x' \Rightarrow \varphi(x) > \varphi(x') $). En esta sección
demostraremos que si esta aplicación existe, $\rightarrow$ es terminante.

\textbf{Fundamento:} Sea $\rightarrow$ terminante. Si no lo fuera, habría una
cadena $x_0 \rightarrow x_1 \rightarrow \dots$, y esta induciría a
$\varphi(x_0) > \varphi(x_1) > \dots$.

\begin{lema}
  Una reducción de ramificación finita $(A,\rightarrow)$ es terminante syss
  existe una aplicación $\varphi : A \rightarrow \mathbb{N}$, que verifique
  $x \rightarrow x' \Rightarrow \varphi(x) > \varphi(x').$
\end{lema}

\begin{demo}
  La implicación hacia la derecha es trivial. Al existir una aplicación
  $\varphi$ que verifique la monotonía, nos da una cadena de $x_i$, y dicha
  cadena es trivialmente terminante.

  Para la demostración de la implicación contraria, supondremos que
  $\varphi(x)$, es el número de sucesores de $x$. Debe ser finito por el lema
  \ref{ramfin}. Como $\rightarrow$ es terminante, y por tanto acíclico (ya que
  no debe tener ningún elemento que $x \xrightarrow{+} x$), para todo
  $x \rightarrow x' \Rightarrow \varphi(x) > \varphi(x')$ (ie, $x'$ tiene menos
  sucesores que $x$).
\end{demo}

%EJEMPLOS

\section{Orden lexicográfico}\label{orlexi}

\begin{defi} 
  Dados dos órdenes estrictos $(A, >_A)$ y $(B, >_B)$, \textbf{el producto
  lexicográfico} $>_{A \times B}$ en $ A \times B$ se define por
  \[(x,y) >_{A \times B} (x',y') :\Leftrightarrow (x >_A x') \vee (x = x' \wedge y >_B y').\]
\end{defi}

\begin{lema}
  El producto lexicográfico de dos órdenes estrictos, es un orden estricto.
\end{lema}

\begin{teor}
  El producto lexicográfico de dos relaciones terminantes, es terminante.
\end{teor}

\begin{demo}
  Por reducción a lo absurdo. Supondremos que $>_A$ y $>_B$ es terminante y que
  la cadena $(a_0,b_0) > (a_1,b_1) > \dots$ es infinita. Como $>_A$ es
  terminante, debe existir un $k$, tal que $a_i = a_{i+1}$ con $i\geq k$. Pero
  esto indica que $>_B$ tiene una cadena infinita, y esto contradice la
  terminación de $>_B$.
\end{demo}

A continuación, extrapolaremos la idea del producto lexicográfico de dos
elementos, a $n$ elementos.

\begin{defi}
  Dado $n$ órdenes estrictos $(A_i,>_i)$, $i=1,\dots,n$ su \textbf{producto
  lexicográfico} $>_{1..n}$ se define por
  \[(x_1,\dots,x_n) >_{1..n}  (y_1,\dots,y_n) :\Leftrightarrow 
    \exists k \geq n. (\forall i < k. x_i = y_i) \wedge x_k >_k y_k.\]
\end{defi} 

\begin{nota} 
  Si todos los $(A_i,>_i)$ son iguales, denotaremos el producto lexicográfico
  como $>^n_{lex}$.
\end{nota}

%EJEMPLO DE VECTORES EN COMPARATIVA CON UN ARBOL

\begin{defi}
  Dado un orden estricto $(A,>)$, el \textbf{orden lexicográfico} $>^*_{lex}$ en $A^*$
  se define por
  \[x >^*_{lex} y :\Leftrightarrow (|x| > |y|) \vee (|x| = |y| \wedge x >^{|x|}_{lex} y).\]
  Donde $|x|$ es la longitud de $x$.
\end{defi} 

\begin{lema}
  Si $>$ es un orden estricto, entonces $>^*_{lex}$ es un orden estricto. Si $>$ es
  terminante, $>^*_{lex}$ también lo es.
\end{lema}

Otra variación del orden lexicográfico será la siguiente.

\begin{defi}
  Diremos que $a >_{Lex} b$ syss $b$ es un prefijo de $a$, ó si
  $a = x \alpha y$, $b = x \beta y$ y $\alpha > \beta$.
\end{defi}

\begin{ejem}
  Si $a>b$, $aaaa >_{Lex} aaa >_{Lex} abba$.
\end{ejem}

\begin{lema}
  Si $>$ es un orden estricto, también lo es $>_{Lex}$.
\end{lema}

\section{Multiconjuntos y su orden}\label{ormult}

Una manera de construir órdenes terminantes es mediante los multiconjuntos. La
definición totalmente informal es que se tratan de conjuntos con elementos
repetidos. Formalmente, los definiremos como sigue,

\begin{defi}
  Un \textbf{multiconjunto} $M$ sobre un conjunto $A$ es una función
  $M : A \rightarrow \mathbb{N}$, tal que $M(x)$ es el número de veces que
  aparece repetido $x$. Un multiconjunto es finito si $\{x \in A : M(x) > 0\}$
  es finito. Mediante $\mathcal{M}(A)$ se denotará el conjunto de todos los
  multiconjuntos finitos sobre $A$.
\end{defi}

Como uno de nuestros principales problemas es la terminación, nos interesan que
los multiconjuntos con los que trabajemos sean finitos.

\begin{ejem} 
  $M = \{a,b,c,b\}$ es un multiconjunto, y es equivalente a
  $\{c,b,b,a\}$. También podemos expresarlo como
  $\{a\mapsto 1, b \mapsto 2, c\mapsto 1\}$
\end{ejem}

\begin{defi} 
  Algunas operaciones básicas y relaciones en $\mathcal{M}(A)$ son, \\
  \begin{tabular}{ll}
    \textbf{Elemento en $M$}
    & $x \in M :\Leftrightarrow M(x) > 0$. \\
    \textbf{Inclusión}     
    & $M \subseteq N :\Leftrightarrow \forall x \in A. M(x) \leq N(x)$. \\ 
    \textbf{Unión}        
    & $(M \cup N)(x) := M(x) + N(x)$.   \\ 
    \textbf{Diferencia}   
    & $(M - N)(x):= M(x) \dotminus N(x)$, \\ 
    & donde $n \dotminus m$ es $n-m$
      si $m \geq n$, 0 en caso contrario.     
  \end{tabular}
\end{defi}

Ya podemos definir un orden para los multiconjuntos a partir de la anterior
definición.

\begin{defi} 
  Dado un orden estricto $>$ en un conjunto A, definimos el \textbf{orden de los
  multiconjuntos} $>_{mul}$ en $\mathcal{M}(A)$ como, $M >_{mul} N$ syss existe
  $X, Y \in \mathcal{M}(A)$, tales que,
  \begin{itemize}
  \item $\emptyset \neq X \subseteq M$,
  \item $N = (M - X) \cup Y$, y 
  \item $\forall y \in Y. \exists x \in X. x>y$.
  \end{itemize}
\end{defi}

Aunque la definición a primera vista puede resultar un poco compleja, lo que
viene a decir es que si $x_1, \dots, x_m$ son elementos de $M$, y
$y_1, \dots, y_m$ son elementos de $N$, $X$ e $Y$ vienen definidos como
\[ \underbrace{x_1, \dots, x_i}_{X} 
   \underbrace{x_{i+1}, \dots, x_m, \dots, y_1, \dots, y_n}_{Y} \] 
y deben cumplir que para todo $y \in Y. \exists x \in X. x>y$. Notase que $X$
e $Y$ se pueden determinar de varias formas.

\begin{lema}
  Si $>$ es un orden estricto, $>_{mul}$ también lo es.
\end{lema}

La propiedad más importante de este capítulo es la siguiente.

\begin{teor}
  $>_{mul}$ es terminante syss $>$ es terminante
\end{teor}

\begin{lema}\label{lemordsyss}
  Si $>$ es un orden estricto y $M,N \in \mathcal{M}(A)$, entonces
  \[ M >_{mul} N\Leftrightarrow M \neq N \wedge 
    \forall n \in N - M. \exists m \in M - N. m > n. \]
\end{lema}

%FALTA HACER LA CONCLUSIÓN

\section{Implementación de los órdenes en Haskell}

En esta sección vamos a implementar en el lenguaje de programación
funcional Haskell, algunas de las definiciones que hemos estudiado en
los anteriores apartados. En el apéndice A se definirán una lista de
funciones básicas de la librería Data.List, en caso de que el lector
no este afianzado con el lenguaje.

Hay que destacar el tipo de dato \texttt{Ordering}, que ya viene predefinido.

\index{\texttt{Ordering}}
\begin{preludio}
data Ordering = EQ | LT | GT
\end{preludio}

Consideramos \texttt{EQ} si $a = b$, \texttt{LT} si $a < b $, y
\texttt{GT} si $a > b$.

\begin{itemize}
\item
\index{\texttt{ordLex}}
    \texttt{(ordLex r xs ys)} es el orden lexicográfico inducido por \texttt{r} sobre \texttt{xs} e \texttt{ys}. 
\begin{sesion}
>>> ordLex compare [1,2] [1,5]
LT
>>> ordLex compare [1,2] [1,2]
EQ
>>> ordLex compare [3,1] [2,4]
GT
>>> ordLex compare [1] [1,3]
LT
>>> ordLex compare [1,3] [1]
GT
>>> let r x y = compare (abs x) (abs y)
>>> ordLex compare [-4,3] [2,5]
LT
> ordLex r       [-4,3] [2,5]
GT
\end{sesion}
        
La definición de la función es,
       
\begin{codigo}
ordLex :: (a -> b -> Ordering) -> [a] -> [b] -> Ordering
ordLex _ [] []  = EQ
ordLex _ [] _   = LT
ordLex _ _  []  = GT
ordLex r (x:xs) (y:ys) = 
    case a of 
      EQ -> ordLex r xs ys
      _  -> a 
    where a = r x y
\end{codigo}      

\item \index{\texttt{ordenMulticonj}} \texttt{(ordenMulticonj r xs
    ys)} es la comparación de los multiconjuntos \texttt{xs} e
  \texttt{ys}. Por ejemplo,
\begin{sesion}
>>> ordenMulticonj (M.fromOccurList [(7,2)])  
  (M.fromOccurList [(7,2)])
EQ
>>> ordenMulticonj (lista2Multiconj [1,2,3,2,5]) 
  (lista2Multiconj [4,5,6,3])
LT
>>> ordenMulticonj (lista2Multiconj [4,5,6,3]) 
  (lista2Multiconj [1,2,3,2,5])
GT
\end{sesion}
        
La definición de la función es
        
\begin{codigo}
ordenMulticonj :: Ord a => M.MultiSet a -> M.MultiSet a 
                  -> Ordering
ordenMulticonj = compare
\end{codigo}

\end{itemize}

\section{Resultados sobre la confluencia}

\begin{defi} 
  Diremos que una relación $\rightarrow$ es \textbf{localmente confluente} si
  \[
    y_1 \leftarrow x \rightarrow y_2 \Rightarrow y_1 \downarrow y_2
  \] 
\end{defi}      

La definición de localmente confluente es mucho mas débil que la de
confluente. No obstante, nos da un resultado importante en relación con la
terminación.

\begin{lema} 
  Una relación terminante es confluente, si es localmente confluente.
\end{lema}

\begin{defi}
  Una relación $\rightarrow$ es \textbf{fuertemente confluente} si,
  \[
    y_1 \leftarrow x \rightarrow y_2 \Rightarrow 
    \exists z.  y_1 \xrightarrow{*} z \xleftarrow{=} y_2
  \] 
\end{defi}

Hay que tener en cuenta, que para $y_1 \leftarrow x \rightarrow y_2$, debe
verificarse tanto $y_1 \xrightarrow{*} z_1 \xleftarrow{=} y_2$, como
$y_1 \xrightarrow{*} z_2 \xleftarrow{=} y_2$

\begin{lema} 
  Una relación fuertemente confluente es confluente.
\end{lema}

\begin{lema} 
  Si $ \ \xrightarrow{*}_1 \ = \ \ \xrightarrow{*}_2$, entonces $\rightarrow_1$
  es confluente syss $\rightarrow_2$ es confluente.
\end{lema}

\begin{lema}
  Si
  $ \ \rightarrow_1 \ \subseteq \ \ \rightarrow_2 \ \subseteq \ \
  \xrightarrow{*}_1 $ entonces $ \ \xrightarrow{*}_1 = \ \ \xrightarrow{*}_2$
\end{lema}

\begin{coro}
  Si
  $ \ \rightarrow_1 \ \subseteq \ \ \rightarrow_2 \ \subseteq \ \
  \xrightarrow{*}_1 $ y $\rightarrow_2 \ $ es fuertemente confluente, entonces
  $\rightarrow_1$ es confluente.
\end{coro}


\begin{defi}
  Una relación $\rightarrow$ tiene la \textbf{propiedad del diamante} si
  \[
    y_1 \leftarrow x \rightarrow y_2 \Rightarrow \exists z. y_1 \rightarrow z \leftarrow y_2. 
  \]
\end{defi} 

\begin{figure}[h]
  \centering
  \begin{tikzpicture}
    \begin{scope}[execute at begin node=$, execute at end node=$]
      \node (v1) at (3,6) {a};
      \node (v2) at (0,3) {b};
      \node (v3) at (6,3) {c};
      \node (v4) at (3,0) {d};
    \end{scope}
        
    \draw [-latex](v2) (v1) edge  node [auto] {} (v2);
    \draw [-latex](v3) (v1) edge  node [auto] {} (v3);
    \draw [-latex](v4) (v2) edge  node [auto] {} (v4);
    \draw [-latex](v4) (v3) edge  node [auto] {} (v4);
  \end{tikzpicture}
        
\caption{Propiedad del diamante}
\end{figure}    

%%% Local Variables:
%%% mode: latex
%%% TeX-master: "SRT_en_Haskell"
%%% End:

%  LocalWords:  ie lex aaaa aaa abba multiconjunto mul List Ordering
%  LocalWords:  EQ LT GT ordLex ys let abs of where ordenMulticonj
%  LocalWords:  fromOccurList Multiconj Ord MultiSet

\chapter{Formalización del concepto de términos}

En este capítulo definiremos el concepto de términos y
sustitución. Este estudio sobre el álgebra universal, es el principal
objeto de estudio de este trabajo. Dedicaremos una sección para cada
definición y propiedades básicas de cada una, incluyendo en ellas su
implementación en Haskell. Dicho código se puede encontrar en
\texttt{Terminos.hs}

\section{Términos}

En esta sección vamos a definir formalmente el concepto intuitivo que
entendemos por término. Por ejemplo, si consideramos $f$ una función de aridad
1, $x$ una variable, entonces $f(x)$ es un término. Para formalizar qué
funciones podemos usar en cada contexto y qué aridad tienen, introducimos el
concepto de signatura.

\begin{defi} 
  Una signatura $\Sigma$ es una familia de conjuntos de símbolos de
  funciones $\{ \Sigma_n \}$, donde para cada símbolo
  $f \in \Sigma_n$, $f$ tiene aridad $n$. Los elementos $\Sigma^{(0)}$
  se llaman símbolos constantes.
\end{defi}

Por ejemplo en álgebra, el concepto de signatura de grupo
es $\Sigma_G = \{e,i,f\}$, donde $e$ es la identidad con aridad 0,
$i$ es la inversa con aridad 1 y $f$ es la suma con aridad 2.

Con esta idea, podemos introducir el concepto de términos.

\begin{defi}
  Sea $\Sigma$ una signatura, y $X$ un conjunto numerable de variables
  tal que, $\Sigma \cap X = \emptyset$. El conjunto de los términos de
  primer orden denotado por $T(\Sigma,X)$ de los $\Sigma$--términos en
  $X$ se define inductivamente como,
  \begin{itemize}
  \item $X \subseteq T(\Sigma,X)$
  \item Si $n \geq 0$, $f \in \Sigma^{(n)}$ y $t_1, \dots, t_n \in
    T(\Sigma, X)$ entonces $f(t_1, \dots, t_n) \in T(\Sigma,X)$ 
  \end{itemize}
\end{defi}

Obsérvese que una variable es un término.  Por ejemplo con el $\Sigma_G$
anterior, $e$, $i$, $f(e,i(e))$ son ejemplos de términos de $\Sigma_G$.

Para su implementación en Haskell necesitamos definir qué
entendemos por una variable. 

\begin{itemize}

\item En primer lugar, definimos los símbolos o nombres como,
\begin{codigo}
type Nombre = String
\end{codigo}

\item Los índices son números naturales.

\begin{codigo}
type Indice = Int
\end{codigo}

\item Los nombres de variables son pares formados por un nombre y un índice.
\begin{codigo}
type Nvariable = (Nombre,Indice)
\end{codigo}
Por ejemplo, $x_3$ se representa como \texttt{("x",3)}.

\end{itemize}

Como un término es una variable o un término compuesto, definimos el
tipo de dato \texttt{Termino}\index{\texttt{Termino}} de la siguiente manera,
\begin{codigo}
data Termino = V Nvariable
             | T String [Termino]
             deriving (Eq, Show)
\end{codigo}

Por ejemplo,
\begin{itemize}
\item $x_2$ se representa por \texttt{(V ("x", 2))}
\item $a$ se representa por \texttt{(T "a" [])}
\item $f(x_2,g(x_2))$ se representa por \texttt{(T "f" [V "x" 2, T "g"
    [V "x" 2]])}
\end{itemize}

A continuación definimos algunos conceptos importantes sobre los
términos.\begin{defi} El conjunto de variables de un término $t$, denotado por
  $V(t)$, se define de manera recursiva como,
  \begin{equation*}
    V(t)=
    \left\lbrace
      \begin{array}{ll}
        \{ t \} & \text{si } t \in X \\
        \displaystyle \bigcup^n_{i=1} V(t_i) & \text{si } t = f(t_1, \dots, t_n)\\
      \end{array}
    \right.
  \end{equation*}
\end{defi}

Por ejemplo el término $f(i(i(e)),f(i(e),x)) \in \Sigma_G$, tiene como
conjunto de variables $\{x\}$.  

Para su implementación en Haskell definiremos la función
\texttt{conjuntoVar t}\index{\texttt{conjuntoVar}}que devuelve la lista de variables del término
\texttt{t}.

\begin{sesion}
ghci> conjuntoVar (V ("x",2))
[("x",2)]
ghci> conjuntoVar (T "f" [T "e" []])
[]
ghci> conjuntoVar (T "f" [T "g" [V ("x",1)], T "h" [], V ("y",1)])
[("x",1),("y",1)]
\end{sesion}

Su definición es,

\begin{codigo}
conjuntoVar :: Termino -> [Nvariable]
conjuntoVar (V x)    = [x]
conjuntoVar (T _ ts) = concatMap conjuntoVar ts
\end{codigo}

\begin{defi} La longitud de un término $t$, denotado por $l(t)$, se
  define de manera recursiva como,
  \begin{equation*}
    l(t)=
    \left\lbrace
      \begin{array}{ll}
        1 & \text{si } t \in X \\
        1 + \displaystyle \sum^n_{i=1} l(t_i) & \text{si } t = f(t_1, \dots, t_n)\\
      \end{array}
    \right.
  \end{equation*}
\end{defi}

Por ejemplo el término $f(i(i(e)),f(i(e),x)) \in \Sigma_G$, tiene
longitud 8.  
En Haskell, creamos la función \texttt{longitudTerm t}\index{\texttt{longitudTerm}} que
calcula la longitud del término \texttt{t}

\begin{sesion}
ghci> longitudTerm (V ("x",2))
1
ghci> longitudTerm (T "f" [T "e" []])
2
ghci> longitudTerm (T "f" [T "g" [V ("x",1)], T "h" [], V ("y",1)])
5
\end{sesion}

Su definición es,

\begin{codigo}
longitudTerm :: Termino -> Int
longitudTerm (V _)    = 1
longitudTerm (T _ xs) = 1 + sum (map longitudTerm xs)
\end{codigo}

\begin{defi} El tamaño de un término $t$, denotado por $\abs{t} $, se
  define de manera recursiva como,
  \begin{equation*}
    \abs{t}=
    \left\lbrace
      \begin{array}{ll}
        0 & \text{si } t \in X \\
        0 & \text{si } t = a \\
        1 + \displaystyle \sum^n_{i=1} \abs{t_i} & \text{si } t = f(t_1, \dots, t_n)\\
      \end{array}
    \right.
  \end{equation*}
\end{defi}

Por ejemplo el término $f(i(i(e)),f(i(e),x)) \in \Sigma_G$, tiene
tamaño 5.  La función que calcula el tamaño de un término \texttt{t},
la llamaremos \texttt{tamanoTerm t}\index{\texttt{tamanoTerm}}.

\begin{sesion}
ghci> tamanoTerm (V ("x",2))
0
ghci> tamanoTerm (T "f" [T "e" []])
1
ghci> tamanoTerm (T "f" [T "g" [V ("x",1)], T "h" [], V ("y",1)])
2
\end{sesion}

Su definición es,

\begin{codigo}
tamanoTerm :: Termino -> Int
tamanoTerm (V _)    = 0
tamanoTerm (T _ []) = 0
tamanoTerm (T _ xs) = 1 + sum (map tamanoTerm xs)
\end{codigo}

Una representación interesante de los términos compuestos es mediante
árboles como se observa en la figura \ref{fig:termino}.

\begin{figure}[h]
  \label{fig:termino}
  \centering
  \begin{tikzpicture}
    \node [circle,draw][label=e]{$f$} child {node
      [circle,draw][label=1]{$i$}
      child{node[circle,draw][label=left:11]{$i$}
        child{node[circle,draw][label=left:111]{$e$}}}} child {node
      [circle,draw][label=2]{$f$} child {node
        [circle,draw][label=21]{$i$}
        child{node[circle,draw][label=right:211]{$x$}}} child {node
        [circle,draw][label=22]{$x$}} };
  \end{tikzpicture}
  \caption{Representación de $f(i(i(e)),f(i(x),x))$ en forma de
    árbol.}
\end{figure}

Las definiciones anteriores se pueden ver más claras con esta
representación. El conjunto de variables de un término $t$ son las
hojas del árbol, la longitud es el número de nodos del árbol, y el
tamaño es el numero de nodos del árbol menos sus hojas. Las variables
que se encuentren en las hojas de este árbol, las llamaremos variables
suelo.

En la figura \ref{fig:termino} cada nodo lleva una numeración. Esto
es lo que conocemos en términos como la posición. Nos resultará útil a
la hora de definir subtérminos y de dotar de cierta estructura al
concepto de término.

\begin{defi}
  Sea $\Sigma$ una signatura, $X$ un conjunto de variables disjuntas
  de $\Sigma$, y $s \in T(\Sigma, X)$. El conjunto posiciones
  $\Pos (s)$ se define de manera recursiva como,
\begin{itemize}

\item Si $s \in X$, entonces $\Pos(s) := {\epsilon}$, donde $\epsilon$
  es la cadena vacía.
\item Si $s = f(s_1, \dots, s_n)$ entonces,
$$ \Pos(s) := \{\epsilon\} \cup \{1p | p \in Pos(s_1) \} \cup \dots \cup \{np | p \in \Pos(s_n) \}$$
\end{itemize}
\end{defi}

El conjunto de $\Pos(s)$ serán cadenas de numeros naturales
positivos. Por ejemplo,
$\Pos f(f(i(i(e)),f(i(x),x))) = \{ \epsilon,1,11,111,2,21,211,22 \}$.

La implementaremos en Haskell bajo el nombre de la función \texttt{conjuntoPos t}\index{\texttt{conjuntoPos}}.

\begin{sesion}
ghci> conjuntoPos (V ("x",1))
[[]]
ghci> conjuntoPos (T "f" [V("x",1)])
[[],[1]]
ghci> conjuntoPos (T "f" [T "i" [T "i" [T "e" []]],
                   T "f" [T "i" [V("x",1)], V("x",1)]])
[[],[1],[1,1],[1,1,1],[2],[2,1],[2,1,1],[2,2]]
\end{sesion}

Su definición es,

\begin{codigo}
conjuntoPos :: Termino -> [[Int]]
conjuntoPos (V _) = [[]]
conjuntoPos (T _ (xs)) = conjuntoPos' (reverse zs)
    where zs = (zip xs [1..])
          conjuntoPos' [] = [[]]
          conjuntoPos' ((t,n):ts) = (conjuntoPos' ts) ++ 
                                    (map (n:) (conjuntoPos t))
\end{codigo}


A continuación, introducimos la noción de subtérmino.
\begin{defi}
  Sea $s \in T(\Sigma, X)$ y $p \in \Pos (s)$. El subtérmino de $s$ en
  la posición $p$, denotado por $s|_p$ se define por inducción
  \begin{equation*}
    \begin{array}{rcl}
      s|_\epsilon & :=  & s \\
      f(s_1,\dots,s_n)|_{iq} & : = & s_i|_q \\
    \end{array}
  \end{equation*}
\end{defi}

De la misma manera, podemos extrapolar este concepto a los árboles. Un
término $t$ es subtérmino del término $s$, si el árbol generado por el
término $t$ es un subárbol para el árbol generado por $s$. 

Implementamos la funcion \texttt{esSubtermino t s} \index{\texttt{esSubtermino}}que comprueba si $t$ es
subtérmino de $s$. Por ejemplo,

\begin{sesion}
ghci> esSubtermino (V ("x",1)) (T "f" [V("x",1)])
True
ghci> esSubtermino (V ("x",2)) (T "f" [V("x",1)])
False
\end{sesion}

Queda implementada de la siguiente manera,

\begin{codigo}
esSubtermino :: Termino -> Termino -> Bool
esSubtermino t s@(T _ xs)
    | t == s = True
    | otherwise = or (map (esSubtermino t) xs)
esSubtermino t s = t == s
\end{codigo}

Una función que usaremos en los capítulos siguientes es
\texttt{ocurre v t}\index{\texttt{ocurre}}, que es la función \texttt{esSubtermino}
restringido a variables. A partir de ahora diremos que $v$ ocurre en
$t$ si la variable $v$ es subtérmino de $t$. Por ejemplo,

\begin{sesion}
ghci> ocurre ("x",2) (V ("x",2)) 
True
ghci> ocurre ("x",3) (V ("x",2)) 
False
ghci> ocurre ("x",2) (T "f" [V ("x",2), T "g" [V ("x",2)]])
True
ghci> ocurre ("y",5) (T "f" [V ("x",2), T "g" [V ("x",2)]])
False
\end{sesion}

Dado la naturaleza de la restriccióna variables, podriamos definir
\texttt{ocurre} como,

\begin{codigo}
ocurre :: Nvariable -> Termino -> Bool
ocurre v t = esSubtermino (V v) t
\end{codigo}

Sin embargo, podemos mejorar la definición sin usar la función
\texttt{esSubtermino},

\begin{codigo}
ocurre :: Nvariable -> Termino -> Bool
ocurre a (V x)    = a == x
ocurre a (T _ ts) = any (ocurre a) ts
\end{codigo}

Con la noción de subtérmino, introducimos una nueva función de
términos en términos. La idea es, a partir de $p\in \Pos(s)$ y de dos
términos $s,t$ sustituir el subtérmino $s|_p$ por $t$. En la
noción de árbol, sería cambiar todo el subárbol $s|_p$ por $t$,
dejando el resto de $s$ intacto. 

% INTRODUCIR UN DIBUJO


Formalmente lo definimos como sigue,

\begin{defi}
Sea $p \in Pos(s)$. Denotamos $s[t]_p$ como,
  \begin{equation*}
    \begin{array}{rcl}
      s[t]_\epsilon & :=  & t \\
      f(s_1,\dots,s_n)[t]|_{iq} & : = & f(s_1,\dots,s_i[t]_q,\dots,s_n) \\
    \end{array}
  \end{equation*}
\end{defi}

La función que calcule el término resultante la llamaremos
\texttt{sustPosSubtermino}\index{\texttt{sustPosSubtermino}}. Debemos tener en cuenta que la de
definición depende de la posición, por ello añadiremos una excepción
si $p$ no esta bien definido. Por ejemplo,

\begin{sesion}
ghci> sustPosSubtermino (T "f" [V ("s",1),V ("s",2),V ("s",3)]) 
                        (V ("t",1)) [2,1]
T "f" [V ("s",1),V ("t",1),V ("s",3)]
sustPosSubtermino (T "f" [V ("s",1),V ("s",2),V ("s",3)]) 
                  (V ("t",1)) []
(V ("t",1))
sustPosSubtermino (V ("s",1)) (V ("t",1)) [2]
*** Exception: No se ha definido bien la lista de posiciones
ghci> sustPosSubtermino (T "f" [V ("s",1),V("s",2),V("s",3)]) 
                        (V ("t",1)) [2,2]
T "f" [V ("s",1),*** Exception: No se ha definido bien la 
                     lista de posiciones
\end{sesion}

Su definición es,

\begin{codigo}
sustPosSubtermino :: Termino -> Termino -> [Int] -> Termino
sustPosSubtermino _ t [] = t
sustPosSubtermino (V _) t [1] = t
sustPosSubtermino (V _) _ _ = error("No se ha definido bien la 
                                     lista de posiciones")
sustPosSubtermino (T f xs) t (i:is) = (T f
    ((take (i-1) xs) ++
    (sustPosSubtermino (xs!!(i-1)) t is) :
    (drop i xs)))
\end{codigo}

\section {Sustituciones}

En esta sección vamos a introducir el concepto de sustituciones entre
términos. De manera general, una sustitución será una aplicación desde
un conjunto de las variables hacia el de los términos. Por ejemplo,
sea $f(x)$ un término $t$ y $\{x \mapsto i(x) \}$ la sustitución
$\tilde{\sigma}$. Si aplicamos $\tilde{\sigma}$ a $t$ obtenemos
$\tilde{\sigma}(t) = f(i(x))$. Formalmente la definiremos como sigue,

\begin{defi} 
  Sea $\Sigma$ una signatura y $X$ un conjunto infinito numerable de
  variables. Una $T(\Sigma, X)$--sustitución, es una función
  $\sigma : X \rightarrow T(\Sigma, X)$ tal que $\sigma(x) \neq x$
  para un número finito de variables $x \in X$. El conjunto de
  variables donde $\sigma(x) \neq x$ se le llama dominio de
  $\sigma$. Lo denotaremos por $\Dom (\sigma)$.
\end{defi}

En definitiva, las sustituciones son de la forma
$$\sigma = \{ x_1 \mapsto \sigma(x_1), \dots, x_n \mapsto \sigma(x_n)\}$$ 
entendiendo que las restantes variables se aplican en ellas mismas.

La implementación del tipo de dato
\texttt{Sustitucion}\index{\texttt{Sustitucion}} en Haskell es,

\begin{codigo}
  type Sustitucion = [(Nvariable,Termino)]
\end{codigo}

Por ejemplo, la sustitución
$\sigma = \{ x_2 \mapsto a, y_4 \mapsto z_5, v_1 \mapsto f(w_2) \}$ 
se representa por 
\begin{sesion}
[(("x",2),T "a" []),
 (("y",4),(V ("z",5))),
 (("v",1),T "f" [V ("w",2)])]
\end{sesion}

Para comprobar si una variable está en el dominio de una sustitución, definimos
la función \texttt{enDominio}\index{\texttt{enDominio}}, tal que \texttt{(enDominio x s)} se verifica si
\texttt{x} está en el dominio de \texttt{s}. Por ejemplo,

\begin{sesion}
ghci> let s = [(("x",2),T "a" []),(("y",4),(V ("z",5)))]
ghci> enDominio ("x",2) s
True
ghci> enDominio ("x",3) s
False
ghci> enDominio ("y",2) s
False
ghci> enDominio ("y",4) s
True
ghci> enDominio ("z",5) s
False
\end{sesion}

Su definición es:

\begin{codigo}
enDominio :: Nvariable -> Sustitucion -> Bool 
enDominio v = any (\(x,_) -> v == x)
\end{codigo}

Con la función \texttt{(aplicaVar s v)}\index{\texttt{aplicaVar}}, calculamos el término
obtenido tras aplicar una sustitución \texttt{s} a una variable
\texttt{v}. Por ejemplo,

\begin{sesion}
let s = [(("x",2),T "a" []),(("y",4),(V ("z",5)))]
ghci> aplicaVar s ("x",2)
T "a" []
ghci> aplicaVar s ("x",3)
V ("x",3)
ghci> aplicaVar s ("y",2)
V ("y",2)
ghci> aplicaVar s ("y",4)
V ("z",5)
ghci> aplicaVar s ("z",5)
V ("z",5)
\end{sesion}

Su definición es:

\begin{codigo}
aplicaVar :: Sustitucion -> Nvariable -> Termino 
aplicaVar [] z = V z 
aplicaVar ((x,y):s) z 
  | x == z    = y
  | otherwise = aplicaVar s z
\end{codigo}

La definifición de sustitución se puede extender para términos
$t \in T(\Sigma, X)$. Sea la extensión $\tilde{\sigma}$ definida como,

\begin{equation*}
  \tilde{\sigma}(t)=
  \left\lbrace
    \begin{array}{ll}
      \tilde{\sigma}(x) = \sigma(x) 
        & \text{si } t \in X \\
      \tilde{\sigma}(f(t_1, \dots, t_n)) = f(\tilde{\sigma}(t_1), \dots, \tilde{\sigma}(t_n)) 
        & \text{si } t = f(t_1, \dots, t_n)
    \end{array}
  \right.
\end{equation*}

A partir de ahora, entenderemos $\sigma$ como una sustitución
extendida, a menos que se indique lo contrario.

Para la computación de esta definición usamos la función
\texttt{(aplicaTerm s t)}\index{\texttt{aplicaTerm}}, que es el término obtenido tras aplicar la
sustitución \texttt{s} al término \texttt{t}. Por ejemplo,

\begin{sesion}
ghci> let s = [(("x",2),T "a" []),(("y",4),(V ("z",5)))]
ghci> aplicaTerm s (T "s" [V ("x",2), T "m" [V ("y",4)]])
T "s" [T "a" [],T "m" [V ("z",5)]]
ghci> aplicaTerm s (T "s" [V ("x",4), T "m" [V ("y",2)]])
T "s" [V ("x",4),T "m" [V ("y",2)]]
\end{sesion}

Su definición es:

\begin{codigo}
aplicaTerm :: Sustitucion -> Termino -> Termino
aplicaTerm s (V x)    = aplicaVar s x
aplicaTerm s (T f ts) = T f (map (aplicaTerm s) ts)
\end{codigo}

Como la sustituciones son muy parecidas a las funciones, podemos
plantearnos la composición entre ellas. Definiéndolas como $\sigma\tau(x) :=
\tilde{\sigma}(\tau(x))$. Veamos que $\sigma\tau(x)$ es una sustitución

\begin{lema}
  La composición de sustituciones es una sustitución.
\end{lema}

\begin{demo}
  $\sigma\tau(x)$ es una función de $X$ a $T(\Sigma, X)$. Como
  $\sigma\tau(x) = x$ ocurre para todas las variables
  $x\in V - (\Dom(\sigma)\cup \Dom(\tau))$, por definición, es una
  sustitución.
\end{demo}

A continuación, introducimos un nuevo concepto imprescindible para la reescritura.

\begin{defi}
  Sea $\Sigma$ una signatura y $X$ un conjunto finito numerable de variables
  disjuntas de $\sigma$. Una identidad es un par
  $(s,t) \in T(\Sigma, X) \times T(\Sigma, X)$. Las denotaremos por
  $s \approx t$, donde $s$ lo denominaremos el lado izquierdo y $t$ el lado
  derecho.
\end{defi} 

La idea de esta definicion es similar a lo que entendemos por
equivalencia. Por ejemplo, en la definición de grupos tenemos la
propiedad asociativa, $(a+b)+c =a+(b+c)$. De la misma manera podemos
respresentar la propiedad como una identidad de $\Sigma_G$,
$f(f(a,b),c) \approx f(a,f(b,c))$.

Podemos relacionar este concepto con lo que hemos visto en el capítulo 1.

\begin{defi}
  Sea $E$ un conjuto de identidades de $\Sigma$. La relación de
  reducción $\rightarrow_E \subseteq T(\Sigma, V) \times T(\Sigma, V)$
  se define como,
  \[
    s \rightarrow_E t \hspace{.3cm} syss \hspace{.3cm} \exists(l,r)
    \in E, p \in \Pos (s), \sigma \text{ una sustitución tal que } s|_p = \sigma(l) \text{ y } t=
    s[\sigma(r)]_p.
  \]
\end{defi}

Volviendo al ejemplo de grupos, su conjunto de identidades es,
\[
  G : = \{f(x,f(y,z)) \approx f(f(x,y),z), f(e,x) \approx x, f(i(x),x)
  \approx e \}
\]
El término $f(i(a), f(a,e))$, con $a$ una constante, tiene una
reducción en la posición $\epsilon$ con la sustitución
$\sigma = \{x \mapsto i(a), y \mapsto a, z \mapsto e \}$, y es
$f(f(i(a),a),e)$. Podemos afirmar nuevamente que
$f(f(i(a),a),e) \rightarrow_G f(e,e)$, en la posición $1$ y la
sustitución $\sigma = \{x \mapsto a \}$.

% Añadir un dibujo


\section{Caracterización de $\xleftrightarrow{*}_E$}

Como hemos visto en el capítulo 1, la clausura transitiva reflexiva de
$\rightarrow_E$ la denotamos por $\xrightarrow{*}_E$, y la clausura
simétrica transitiva reflexiva por $\xleftrightarrow{*}_E$. En la
última parte de este capítulo, vamos a dar una caracterización de
$\xleftrightarrow{*}_E$ que usaremos en capítulos posteriores.

\begin{defi} 
  Sea $\equiv$ una relación binaria en $T(\Sigma,V)$
  \begin{enumerate}
  \item La relación $\equiv$ es cerrada bajo sustituciones syss $s \equiv t$
    implica $\sigma(s) \equiv \sigma(t)$, para todo $s,t$ y $\sigma$. 
  \item La relación $\equiv$ es cerrada bajo $\Sigma$--operaciones syss $s_1
    \equiv t_1, \dots, s_n \equiv t_n$ implica $f(s_1, \dots, s_n) \equiv
    f(t_1, \dots, t_n)$, para todo $n \geq 0, f, s_1, \dots s_n, t_1, \dots
    t_n.$ 
  \item La relación $\equiv$ es compatible con $\Sigma$--operaciones syss $ s
    \equiv t$ implica $f(s_1, \dots, s_{i-1},s, s_{i+1},\dots, s_n) \equiv
    f(s_1, \dots, s_{i-1},t, s_{i+1},\dots, s_n)$ para todo $n \geq 0, f, i =
    1, \dots, n,$ y $s_1, \dots, s_{i-1},s,t, s_{i+1},\dots, s_n$. 
  \item La relación $\equiv$ es compatible con el $\Sigma$--contexto syss $s
    \equiv s'$ implica $t[s]_p \equiv t[s']_p$ para todo $t$ y $p$. 
  \end{enumerate}       
\end{defi}

Una consecuencia directa a partir de la definición de $\rightarrow_E$ es,

\begin{lema} \label{lema:2.1}
  Sea $E$ un conjunto de $\Sigma$--identidades. La relación de reducción de
  $\rightarrow_E$ es cerrada bajo sustitución y compatible con las
  $\Sigma$--operaciones.
\end{lema}

\begin{demo}
  Es trivial pues si sabemos que $s_i \equiv t_i \forall i$, entonces
  $\sigma(s_i) \equiv \sigma(t_i) \forall i$, y
  $f(s_1, \dots, s_n) \equiv f(t_1, \dots, t_n)$ por la propia
  definición de relacción de reducción.
\end{demo}

Otro lema importante para la caracterización es,

\begin{lema} \label{lema:2.2}
  Sea $\equiv$ una relación binaria en $T(\Sigma,V)$.
  \begin{enumerate}
  \item La relación $\equiv$ es compatible con las
    $\Sigma$--operaciones syss es compatible con el
    $\Sigma$--contexto.
  \item Si $\equiv$ es reflexiva y transitiva, entonces es compatible
    con las $\Sigma$--operaciones syss es cerrado bajo las
    $\Sigma$--operaciones.
  \end{enumerate}
\end{lema}

\begin{demo}
  La primera propiedad de derecha a izquierda es trivial por la
  definición de posición. De izquerda a derecha se realiza mediante
  una inducción de $p$, como para $p= \epsilon$ la afirmación es
  cierta, también lo será para $p=p_1$, por tanto también para
  $p=p_1 p_2 \dots$.  La segunda propiedad de izquierda a derecha es
  cierta por la transitividad de $\equiv$. De derecha a izquierda se
  usa que la relación es reflexiva.
\end{demo}

\begin{teor} \label{teor:2.1}
  Sea $E$ un conjunto de $\Sigma$--identidades. La relación
  $\xleftrightarrow{*}_E$ es la relación de equivalencia más pequeña en
  $T(\Sigma,V)$ que contiene $E$ que es cerrada bajo sustituciones y
  $\Sigma$-operaciones.
\end{teor}

\begin{demo}
  $\xleftrightarrow{*}_E$ es una relación de equivalencia reflexiva y
  transitiva. Con el lema \ref{lema:2.1} tras aplicar inducción, se puede
  ver que es cerrada bajo sustituciones y compatible con
  $\Sigma$--operaciones. Por el lema \ref{lema:2.2}, también es cerrada bajo
  $\Sigma$--operaciones.

  Vamos a suponer que existe $\equiv$ relación de equivalencia en
  $T(\Sigma,v)$ que contiene a $E$, y es cerrada bajo sustituciones y
  $\Sigma$--operaciones. Nuestro objetivo es comprobar que
  $\xleftrightarrow{*} \subseteq \equiv$, para ello, demostraremos que
  si $s \rightarrow_E t$ entonces $s \equiv t$.

  Partiendo de $s \rightarrow_E$, por la definición de relación de
  reducción, existe $(l,r) \in E$, $p$ posición de $s$, y $\sigma$ una
  sustitución tal que $s|_p = \sigma(l)$ y $t= s [\sigma(r)]_p$. 

  Como $\equiv$ contiene a $E$ entonces $l \equiv r$, por tanto,
  $\sigma(l) \equiv \sigma(r)$ por ser cerrada bajo sustituciones.  Al
  ser $\equiv$ reflexiva y cerrada bajo $\Sigma$--operaciones y
  compatible con $\Sigma$--operaciones, por el lema \ref{lema:2.2} también
  es compatible con $\Sigma$--contexto y por la definición de la
  propiedad, $s = s[\sigma(l)]|_p \equiv s[\sigma(r)]|_p = t$.
\end{demo}

La demostración del teorema \ref{teor:2.1}, nos indica una manera de
calcular $\xleftrightarrow{*}_E$ a partir de $E$ mediante las
relaciones de reflexividad, simetría, transitividad, sustituciones,
y $\Sigma$--operaciones. La aplicacion de estas reglas en este
contexto, es lo que llamamos lógica ecuacional. Las reglas se pueden escribir así,

\[
  \begin{array}{lc}
    $Reflexiva$     & \dfrac{}{E \vdash t \approx t} \\ \\
    $Sim\'{e}trica$         & \dfrac{E \vdash s \approx t}{E \vdash t \approx s} \\ \\
    $Transitiva$    & \dfrac{E \vdash s \approx t  \hspace{.8cm}   E \vdash t \approx u}{E \vdash s \approx u} \\ \\
    $Sustitución$   & \dfrac{E \vdash s\approx t}{E \vdash \sigma(s)\approx \sigma(t)} \\ \\
    \Sigma$--operaciones$    & \dfrac{E \vdash s_1\approx t_1 \hspace{.8cm} \dots \hspace{.8cm} E \vdash s_n\approx t_n}{E \vdash f(s_1, \dots, s_n) \approx f(t_1,\dots,t_n)}
  \end{array} 
\]


%%% Local Variables:
%%% mode: latex
%%% TeX-master: "SRT_en_Haskell"
%%% End:

\chapter{El problema de la unificación y normalización de términos}

En este capítulo vamos definir los problemas de decisión en el
razonamiento ecuacional. Nos centraremos en los sistemas de
reescritura de términos y la unificación.  Para ello consideraremos la
siguiente definición.

\begin{defi}
  Dado un conjunto de identidades $E$ y una ecuación $s \approx t$, diremos que la ecuación es, \\
  $\begin{array}{llll}
    $\textbf{Válida}$        &  $ en $E$ syss s $s \approx_E t  \\
    $\textbf{Satisfacible}$  &  $ en $E$ syss hay una sustitución $\sigma$ tal que $\sigma(s) \approx_E \sigma(t) \\
   \end{array} 
  $
\end{defi}

En la siguiente sección obtendremos varias propiedades interesantes sobre $\approx_E$.

\section{Estudio de $\approx_E$}

%FALTA REFERENCIA AQUI
En el teorema 1.2, vimos que $\xleftrightarrow{*}$ era
decidible si $\rightarrow_E$ si es convergente. De esta manera
podíamos computar $\downarrow_E$. Por ello es importante saber si
podemos decidir si un término implica a otro término. Esto se llama el
problema del emparejamiento, dado dos términos $s$ y $l$, ver si
existe una sustitución $\sigma(l) = s$, y computar $\sigma$ si
existe. 

Comprobaremos que el problema del emparejamiento y de la
satisfacibilidad están muy relacionados.

\begin{teor} \label{teor:3.1}
  Si $E$ es finito y $\rightarrow_E$ es covergente, entonces
  $\approx_E$ es decidible.
\end{teor}

\begin{demo}
%REFERENCIA
  Por el teorema 1.2, sabemos que $s \xleftrightarrow{*}_E t$ syss
  $s\downarrow_E = t\downarrow_E$.  La forma normal del operador
  $\downarrow_E$ es computable, pues podemos decidir si un termino $u$
  es forma normal, y si no lo fuese, podemos computar un $u'$ tal que
  $u \rightarrow_E u'$, y $u'$ fuese una forma normal.
	
  Para ver si $u$ es una forma normal, tenemos que comprobar para
  todas las identidades de $E$ y todas la posiciones $p \in Pos(u)$ si
  hay una sustitución $\sigma$ tal que $u|_p = \sigma l$. Como el
  problema del emparejamiento es decidible y $E$ es finito, $u$ es una
  forma normal o podemos reducir $u$ en un proceso iterativo finito a
  una formal normal. Este proceso termina pues $\rightarrow_E$ es
  terminante.
\end{demo}

Por este resultado, surge un interés en las reducciones convergentes.
El problema de la palabra para $E$ trata de decidir $s \approx t$ para
arbitrarios $s,t, \in T(\Sigma, V)$. También tenemos el problema de
las palabras básicas para $E$, que no es más que el problema de la
palabra restringido a términos básicos $s,t$.

\begin{ejem}
  Considerando la signatura $\Sigma = \{*, S, K\}$ donde $*$ es una
  función binaria y $S$ y $K$ son constantes. El conjunto de
  identidades considerado es,

  \[E := \{((S \cdot x) \cdot y) \cdot z \approx (x \cdot z) \cdot (y
    \cdot z),(K \cdot x) \cdot y \approx x \} \]

  Como se puede dar cualquier término a partir de uno básico, esto
  implica que el problema de la palabra para E sea indecidible.
\end{ejem}

\section{Sistemas de reescritura de terminos}

En esta sección vamos a introducir la definición de regla y sistema de
reescritura, así como una propiedad.

\begin{defi}
  Una regla de reescritura es una identidad $l \approx r$ tal que $l$
  no es una variable y $\Var(l) \supseteq \Var(r)$. En ese caso,
  escribiremos $l \rightarrow r$. Un sistema de reescritura de
  términos (SRT) es un conjunto de reglas de escritura. Una expresión
  reducible es una instancia del lado derecho de una regla de
  reescritura. Contraer la expresión reducible significa sustituirla
  por la instancia del lado derecho de la regla.
\end{defi}

Como cualquier SRT es un conjunto de identidades, $\rightarrow_R$ está
bien definido.

A partir del teorema \ref{teor:3.1}, deducimos el siguiente resultado,

\begin{teor}
  Si $R$ es un SRT convergente finito, $\approx_R$ es decidible: $s
  \approx_R t$ syss $s \downarrow_R = t \downarrow_R$.
\end{teor}

Veamos ahora una de las propiedades básicas de $\rightarrow_R$.

\begin{defi}
  Una relación en $T(\Sigma, V)$ es una relación de reescritura syss
  es compatible con las $\Sigma$-operaciones y cerrada bajo
  sustitución.
\end{defi}

Por el \ref{lema:2.2}, sabemos que $\rightarrow_R$ es una relación de
reescritura. Por inducción de las reglas de derivación,
$\xrightarrow{*}_R$ y $\xrightarrow{+}_R$ también son relaciones de
reescritura. Para $\xleftrightarrow{*}_R$ lo demuestra el teorema \ref{teor:2.1}

\section{Unificación sintática}

En esta sección nos centraremos en estudiar la unificación. La
unificación es el proceso de resolver la satisfacibilidad de un
problema. Dado $E$, $s$ y $t$, hay que encontrar una sustitución
$\sigma$ tal que $\sigma s \approx _E \sigma t$. Si $s$ y $t$ son
básicos, el problema pasaría a ser un problema de palabras básicas. En
esta sección estudiaremos el caso especial $E = \emptyset$ que es
importante para algunos algorítmos simbólicos de computación. El que
especialmente nos importa es un algoritmo que comprueba la confluencia
a partir de pares críticos, que veremos mas adelante.

La sustitución o el unificador que busquemos va a
depender del caso en el que nos encontremos. Por ejemplo $x \approx^?
f(y)$ tiene varios unificadores $\{x \mapsto f(y) \}$, $\{x \mapsto
f(a), y \mapsto a\}$. No siempre puede existir uno o varios
unificadores por ejemplo, $f(x) \approx^? g(x)$ no tiene
unificador. Para el caso en el que existan varios unificadores, nos
gustaría encontrar el más general,

\begin{defi}
  Una sustitución $\sigma$ es mas general que una sustitución
  $\sigma'$ si hay una sustitución $\delta$ tal que $ \sigma' = \delta
  \sigma$. En este caso escribimos $\sigma \lesssim \sigma'$. Diremos
  que $\sigma'$ es una instancia de $\sigma$.
\end{defi}

En el ejemplo anterior, si $\sigma = \{x \mapsto f(y) \}$ y
$\sigma' = x \mapsto f(a), y \mapsto a\}$ entonces
$\sigma \lesssim \sigma'$ porque $\sigma' = \delta \sigma$ donde
$\delta = \{ y \mapsto a \}$.

Una propiedad inmediata de la definición es,
%escribir intr de \sim
\begin{lema} \label{lema:3.1}
  $\sigma \sim \sigma '$ syss existe una sustitución inyectiva $p$ que
  sea biyectiva para las variables tal que $\sigma = p \sigma '$
\end{lema}

Como nuestro objetivo es resolver conjuntos de ecuaciones,
formalizaremos el problema de la unificación con la siguiente definición,

\begin{defi}
  Un \textbf{unificador} de un conjunto finito de ecuaciones
  $S = \{ s_1 =^? t_1, \dots, s_n =^? t_n\}$ es una sustitución $\sigma$
  tal que $\sigma s_i = \sigma t_i$ para $i = 1, \dots, n.$ Por $U(S)$
  se denota el conjunto de todos los unificadores de $S$. Se dice que
  $S$ es \textbf{unificable} si $U(S) \neq \emptyset$. El
  \textbf{problema de la unificación} de $S$ consiste en calcular sus
  unificadores.
\end{defi}

Otra definición importante es,

\begin{defi} 
  Una sustitución $\sigma$ es el \textbf{unificador de máxima
    generalidad} (UMG) de $S$ si $\sigma$ es el menor elemento de
  $U(S)$; es decir,
  \begin{enumerate}
  \item $\sigma \in U(s)$
  \item $\forall \sigma' \in U(s) \sigma \lesssim \sigma'$
  \end{enumerate}
\end{defi}

Por ejemplo, $\sigma ' := \{ x \mapsto z, y \mapsto z \}$ es un
unificador de $x=f =^? y$, pero no es UMG, pues es menos general que
$\sigma = {x \mapsto y}$.

\begin{defi}
  Una sustitución $\sigma$ es \textbf{idempotente} si $\sigma = \sigma \sigma$.
\end{defi}

\begin{lema} \label{lema:3.2} Una sustitución $\sigma$ es indempotente
  syss $\Dom(\sigma) \cap \VRan(\sigma) = \emptyset $
\end{lema}

El ultimo resultado de este capítulo se usará en la implementación del
problema de la unificación.

\begin{teor}
  Si un problema de unificación $S$ tiene solución, entonces tiene un
  UMG idempotente.
\end{teor}

\begin{demo} %volver a mirar esta demo
  El lema \ref{lema:3.1} nos dice que los UMG son unicos respecto a
  una sustitución $p$ inyectiva y biyectiva para las variables. Por
  ello, incluso para los idempotentes UMGs no serán únicos.
\end{demo}

\section{Unificación por transformación}

La unificación puede ser representada como la aplicación iterada de
transformaciones en un conjunto de ecuaciones. Por ejemplo el
algoritmo que seguimos para resolver sistemas lineales,

% HACER EJEMPLO

\begin{defi}
  El problema de la unificación de 
  $S = \{ x_1 =^? t_1, \dots, x_n =^? t_n \}$ 
  está en su \textbf{forma resuelta}, si las variables $x_i$ son distintas
  entre sí y ninguna $x_i$ ocurre en $t_i$. En este caso definimos, 
  $\vec{S} := \{x_1 \mapsto t_1, \dots, x_n \mapsto t_n \}$.
\end{defi}

Varias propiedades básicas de $\vec{S}$ son,

\begin{lema} \label{lema:3.3}
  Si $S$ está en un su forma resuelta entonces, 
  $\sigma = \sigma \vec{S}$ 
  para todo $\sigma \in U(S)$.
\end{lema}

\begin{demo}
  Sea $S = \{ x_1 =^? t_1, \dots, x_n =^? t_n \}$. Probaremos por
  casos que $\forall x \in V. \sigma x = \sigma \vec{S} x$

  Para el caso $x \in \{x_1, \dots, x_n \}$, supongamos sin perdida de
  generalidad que $x = x_m$. Entonces
  $\sigma x = \sigma t_k = \sigma \vec{S} x$ pues $\sigma \in U(S)$.

  Para el caso $x \not \in \{x_1, \dots, x_n \}$. Entonces
  $\sigma x = \sigma \vec{S} x$, pues $\vec{S} x = x$.
\end{demo}

\begin{lema}
  Si $S$ está en su forma resuelta entonces $\vec{S}$ es un UMG
  idempotente de $S$.
\end{lema}

\begin{demo}
  Usando el lema \ref{lema:3.2}, deducimos que ninguno de los $x_i$
  ocurren en $t_i$. Por la misma razón tenemos que
  $\vec{S} x_i = t_i = \vec{S} t_i$, es decir $\vec{S} \in U(S)$.
  Finalmente $\vec{S}$ es un UMG por el lema \ref{lema:3.3},
  $\vec{S} \lesssim \sigma$ para todo $\sigma \in U(S)$.
\end{demo}

Es decir, una vez que tengamos la forma resuelta de $S$, podemos extraer
un UMG idempotente.

Vamos a dar una serie de reglas de transformación para conseguirlo.
\\
$
\begin{array}{ll}
  $Borrar$       & $Elimina ecuaciones triviales$ \\ 
  $Descomponer$  & $Reemplaza ecuaciones entre términos por ecuaciones entre $ \\
                 & $sus subtérminos$ \\ 
  $Orientar$     & $Mueve variables hacia el lado izquierdo$ \\ 
  $Resuelve$     & $Afianza las soluciones, eliminando la variable resuelta en el$ \\
                 & $resto del problema$
\end{array} 
$
\\
\\
Que formalmente significan,
\\
\\
$
\begin{array}{llll}
  $Borrar$      
    & \{t =^? t  \} \sqcup S                             
    & \Longrightarrow & S \\ 
  $Descomponer$ 
    & \{f(\overline{t_n}) =^? f(\overline{u_n}) \} \sqcup S  
    & \Longrightarrow & \{t_1 =^? u_1, \dots, t_n =^? u_n\} \cup S\\ 
  $Orientar$    
    & \{t =^? x \} \sqcup S                              
    & \Longrightarrow & \{x=^? t \} \cup S$ si $t \not\in V \\ 
  $Resuelve$    
    & \{x =^? t \} \sqcup S                              
    & \Longrightarrow & \{x=^? t \} \cup \{x \mapsto t \}(S) \\
    &                                                    
    &                 & $si $ x \in \Var(S) - \Var(t) \\
\end{array} 
\\
$ Donde $\sqcup$ denota la unión disjunta.

%Ejemplo

\begin{defi} 
  La función $\Unify(S)$ devuelve un unificador de máxima generalidad si
  existe, y falla si no. Tiene la siguiente definición, \\ 
  $\begin{array}{lll}
    $\textbf{\Unify(S)}$      
    & = & $Mientras exista un $T$ tal que $S \Longrightarrow T, S:= T$.$ \\
    & & $Si $S$ es una forma resuelta entonces devuelve $\vec{S}$, si no, falla.$ \\
   \end{array} 
  $
\end{defi}

Este algoritmo es no determinista, pues siempre se puede aplicar más
de una regla de transformación. Por ello, el algoritmo escogerá una
regla arbitraria. La terminación de $\Unify$ dependerá de la
terminación de $\Longrightarrow$.

\begin{lema} 
  $\Unify$ termina para todos las entradas.
\end{lema}

\begin{lema}
  Si $S \Longrightarrow T$ entonces $U(S) = U(T)$
\end{lema}

\begin{lema} 
  Si \Unify(S) devuelve una sustitución $\sigma$, $\sigma$ es un UMG
  idempotente de $S$
\end{lema}

% La prueba de S es resoluble, necesita de los dos siguientes
% enunciados

\begin{lema}
  Una ecuación $f(\sigma(\overline{s_m})) =^?  g(\overline{t_n})$
  donde $f \not = g$, no tiene solución.
\end{lema}

\begin{lema}
  Una ecuación $x =^ t$, donde $x \in \Var(t)$ y $x \not = t$, no
  tiene solución.
\end{lema}

\begin{lema} 
  Si $S$ es resoluble, $\Unify(S)$ no falla.
\end{lema}

% HAY UN PARRAFO AQUI, PERO CREO QUE NO SE USARA MAS ADELANTE
\section{Implementación de la unificación y la reescritura de términos
  en Haskell}

En esta sección vamos a implementar en el lenguaje de programación
funcional Haskell los algorítmos de este capítulo. Usaremos algunas
funciones del capítulo anterior, que se encuentran en la libreria
\texttt{Terminos.hs}.

Empezamos implementando los errores. Estos son causados por
\texttt{UNIFICACION} (si el sistema no tiene solución) o
\texttt{REGLA} (si durante la reescritura no se puede aplicar ninguna
regla más).
\begin{codigo}
data ERROR = UNIFICACION
           | REGLA
           deriving (Eq, Show)
\end{codigo}

Una ecuación es un par de términos
\begin{codigo}
type Ecuacion = (Termino,Termino)
\end{codigo}

Un sistema de ecuaciones es una lista de ecuaciones
\begin{codigo}
type Sistema = [Ecuacion]
\end{codigo}

\subsection{Algorítmo de unificiación}

\index{\texttt{(unificacion t1 t2)}} \texttt{(unificacion t1 t2)} es
\begin{itemize}
\item \texttt{(Right s)} si los términos \texttt{t1} y \texttt{t2} son
  unificables y \texttt{s} es un unificador del máxima generalidad de
  \texttt{t1} y \texttt{t2}.
\item \texttt{(Left UNIFICACION)} si \texttt{t1} y \texttt{t2} no son
  unificables.
\end{itemize}
Por ejemplo,
\begin{sesion}
ghci> unificacion (T "f" [V ("x",0), T "g" [V ("z",0)]])
                  (T "f" [T "g" [V ("y",0)], V ("x",0)])
Right [(("z",0),V ("y",0)),(("x",0),T "g" [V ("y",0)])]
ghci> unificacion (T "f" [V ("x",0),T "b" []])
                  (T "f" [T "a" [],V ("y",0)])
Right [(("y",0),T "b" []),(("x",0),T "a" [])]
ghci> unificacion (T "f" [V ("x",0),V ("x",0)])
                  (T "f" [T "a" [],T "b" []])
Left UNIFICACION
ghci> unificacion (T "f" [V ("x",0),T "g" [V ("y",0)]])
                  (T "f" [V ("y",0),V ("x",0)])
Left UNIFICACION
\end{sesion}

Su código es,
\begin{codigo}
unificacion :: Termino -> Termino -> Either ERROR Sustitucion
unificacion t1 t2 = unificacionS [(t1,t2)] []
\end{codigo}
 


\index{\texttt{(unificacionS es s)}} \texttt{(unificacionS es s)} es
\begin{itemize}
\item \texttt{(Right s')} si el sistema \texttt{s(es)}, obtenido
  aplicando la sustitución \texttt{s} a la ecuaciones de \texttt{es},
  es unificable y \texttt{s'} es un unificador del máxima generalidad
  de \texttt{s(es)}.
\item \texttt{(Left UNIFICACION)} si \texttt{s(es)} no es unificable.
\end{itemize}
Por ejemplo,
\begin{sesion}
ghci> unificacionS [(T "f" [V ("x",0), T "g" [V ("z",0)]],
                   T "f" [T "g" [V ("y",0)], V ("x",0)])]
                   []
Right [(("z",0),V ("y",0)),(("x",0),T "g" [V ("y",0)])]
ghci> unificacionS [(T "f" [V ("x",0),T "b" []],
                   T "f" [T "a" [],V ("y",0)])]
                   []
Right [(("y",0),T "b" []),(("x",0),T "a" [])]
ghci> unificacionS [(T "f" [V ("x",0),V ("x",0)],
                   T "f" [T "a" [],T "b" []])]
                   []
Left UNIFICACION
ghci> unificacionS [(T "f" [V ("x",0),T "g" [V ("y",0)]],
                   T "f" [V ("y",0),V ("x",0)])]
                   []
Left UNIFICACION
\end{sesion}

Su código es,
\begin{codigo}
unificacionS :: Sistema -> Sustitucion -> Either ERROR Sustitucion
unificacionS [] s = Right s
unificacionS ((V x,t):ts) s 
  | V x == t  = unificacionS ts s
  | otherwise = reglaElimina x t ts s
unificacionS ((t,V x):ts) s = 
  reglaElimina x t ts s
unificacionS ((T f ts1, T g ts2):ts) s 
  | f == g    = unificacionS (zip ts1 ts2 ++ ts) s
  | otherwise = Left UNIFICACION
\end{codigo}

\index{\texttt{(reglaElimina x t es s)}} \texttt{(reglaElimina x t es s)} es
\begin{itemize}
\item \texttt{(Left UNIFICACION)}, si \texttt{x} ocurre en \texttt{t}.
\item \texttt{(Right s')} es la solución del sistema obtenido al
  aplicarle a \texttt{es} la sustitución [x/t] en el entorno obtenido
  componiendo la sustitución \texttt{s} con [x/t], en caso
  contrario.
\end{itemize}
Por ejemplo,
\begin{sesion}
ghci> reglaElimina ("x",0) (T "f" [V ("x",0)]) [] []
Left UNIFICACION
ghci> reglaElimina ("x",0) (T "f" [V ("x",1)]) [] []
Right [(("x",0),T "f" [V ("x",1)])]
ghci> reglaElimina ("x",0) (T "f" [V ("x",1)]) [] [(("y",2),V ("x",0))]
Right [(("x",0),T "f" [V ("x",1)]),(("y",2),T "f" [V ("x",1)])]
ghci> reglaElimina ("x",0) (T "f" [V ("x",1)]) [(V ("x",1),V ("x",0))] []
Left UNIFICACION
ghci> reglaElimina ("x",0) (T "f" [V ("x",1)]) [(V ("x",1),V ("y",0))] []
Right [(("x",1),V ("y",0)),(("x",0),T "f" [V ("y",0)])]
ghci> reglaElimina ("x",0) (T "f" [V ("x",1)]) [(V ("y",1),V ("x",0))] []
Right [(("y",1),T "f" [V ("x",1)]),(("x",0),T "f" [V ("x",1)])]
\end{sesion}

Su código es,
\begin{codigo}
reglaElimina :: Nvariable -> Termino -> Sistema -> Sustitucion 
             -> Either ERROR Sustitucion
reglaElimina x t es s 
  | ocurre x t = Left UNIFICACION
  | otherwise  = unificacionS es' s'
  where es' = [(aplicaTerm [(x,t)] t1, aplicaTerm [(x,t)] t2) 
                 | (t1,t2) <- es] 
        s'  = (x,t) : map (\(y,u) -> (y, aplicaTerm [(x,t)] u)) s
\end{codigo}
 
%¿Incluyo a modo de ejemplo lo obtenido con el Debug.Trace?

\subsection{Equiparación de términos}

\index{\texttt{(equiparacion t1 t2)}} \texttt{(equiparacion t1 t2)} es
\begin{itemize}
\item \texttt{(Right s)} si \texttt{t1} y es una instancia de \texttt{t1} y \texttt{s} es un equiparador
  del máxima generalidad de \texttt{t1} con \texttt{t2}.
\item \texttt{(Left UNIFICACION)} en caso contrario.
\end{itemize}
Por ejemplo,
\begin{sesion}
ghci> let t1 = T "f" [V ("x",0), T "a" []]
ghci> let t2 = T "f" [T "b" [],  T "a" []]
ghci> equiparacion t1 t2
Right [(("x",0),T "b" [])]
ghci> equiparacion t2 t1
Left UNIFICACION
ghci> let t3 = T "f" [V ("x",0), V ("y",0)]
ghci> equiparacion t1 t3
Left UNIFICACION
\end{sesion}

Su código es,
\begin{codigo}
equiparacion :: Termino -> Termino -> Either ERROR Sustitucion
equiparacion t1 t2 = equiparacionS [(t1,t2)] []
\end{codigo}

\index{\texttt{(equiparacionS es s)}} \texttt{(equiparacionS es s)} es
\begin{itemize}
\item \texttt{(Right s')} si la derecha de las ecuaciones de
  \texttt{s(es)} y es una instancia de la izquierda y \texttt{s'} es
  un equiparador de máxima generalidad de la izquierda con la derecha
\item \texttt{(Left UNIFICACION)} en caso contrario.
\end{itemize}
Por ejemplo,
\begin{sesion}
ghci> let t1 = T "f" [V ("x",0), T "a" []]
ghci> let t2 = T "f" [T "b" [],  T "a" []]
ghci> equiparacionS [(t1,t2)] []
Right [(("x",0),T "b" [])]
ghci> equiparacionS [(t2,t1)] []
Left UNIFICACION
ghci> let t3 = T "f" [V ("x",0), V ("y",0)]
ghci> equiparacionS [(t1,t3)] []
Left UNIFICACION
\end{sesion}

Su código es,
\begin{codigo}
equiparacionS :: Sistema -> Sustitucion -> Either ERROR Sustitucion
equiparacionS [] s = Right s
equiparacionS ((V x,t):es) s 
    | not (enDominio x s)   = equiparacionS es ((x,t):s)
    | aplicaVar s x == t    = equiparacionS es s
    | otherwise             = Left UNIFICACION
equiparacionS ((_,V _):_) _ = Left UNIFICACION
equiparacionS ((T f ts1, T g ts2):es) s 
    | f == g    = equiparacionS (zip ts1 ts2 ++ es) s
    | otherwise = Left UNIFICACION
\end{codigo}

\subsection{Reescritura de términos}

\index{\texttt{(reescribe es t)}} \texttt{(reescribe es t)} es
\begin{itemize}
\item \texttt{(Right s')}, donde \texttt{s'} es el término obtenido
  reescribiendo \texttt{t} con la primera regla de es que con la que
  se pueda reescribir.
\item \texttt{(Left REGLA)} en caso contrario.
\end{itemize}
Por ejemplo,
\begin{sesion}
ghci> reescribe [(V ("x",1), V ("y",1)),
                 (V ("x",2), V ("y",2))]
                 (V("x",2))
Right (V ("y",2))
ghci> reescribe [(V ("x",1), V ("y",1)),
                 (V ("x",1), V ("y",2))]
                 (V("x",2))
Left REGLA
ghci> reescribe [(V ("x",2),T "f" [V ("x",3)]),
                 (V ("x",3),T "f" [V ("x",4)])]
                 (V("x",2))
Right (T "f" [V ("x",3)])
ghci> reescribe [(V ("x",1),T "a" [])]
                 (T "f" [V ("x",1), T "b" []])
Right (T "f" [T "a" [],T "b" []])
\end{sesion}

Su código es,
\begin{codigo}
reescribe :: Sistema -> Termino -> Either ERROR Termino
reescribe [] _ = Left REGLA
reescribe ((l,r):es) t
    | a == Left UNIFICACION = reescribe es t
    | b == t = reescribe es t
    | otherwise = Right b
    where a = equiparacion l r
          b = aplicaTerm (elimR a) t
          elimR (Right x) = x
\end{codigo}
 
\index{\texttt{(formaNormal es t)}} \texttt{(formaNormal es t)} es la forma normal de \texttt{t} respecto
de \texttt{es}. Por ejemplo,

\begin{sesion}
ghci> formaNormal [(V ("x",1), V ("y",1)),
                   (V ("x",2), V ("y",2))]
                   (V("x",2))
V ("y",2)
ghci> formaNormal [(V ("x",1), V ("y",1)),
                   (V ("x",2), V ("y",2))]
                   (V("x",3))
V ("x",3)
ghci> formaNormal [(V ("z",1), V ("a",1)),
                   (V ("x",1), T "g" [V ("z",1)])]
                   (T "f" [V ("x",1)])
T "f" [T "g" [V ("a",1)]]
\end{sesion}

Su código es,
\begin{codigo}
formaNormal :: Sistema -> Termino -> Termino
formaNormal es (V x)
  | isRight a = elimR a
  | otherwise = V x
  where a = reescribe es (V x)
        elimR (Right r) = r
formaNormal es (T f ts)
  | a == Left REGLA = u
  | otherwise       = formaNormal es (elimR a)
  where u = T f (map (formaNormal es) ts)
        a = reescribe es u
        elimR (Right r) = r
\end{codigo}






 
%%% Local Variables:
%%% mode: latex
%%% TeX-master: "SRT_en_Haskell"
%%% End:

\chapter{Terminación}

Como hemos podido comprobar en los anteriores capítulos, es importante
que nuestros sistemas de reescritura tengan la propiedad de la
terminación. Sin embargo, en la primera sección de este capítulo
demostraremos que el problema de decidir si un sistema de reescritura
es terminante, es indecidible. Aunque en casos mas restringidos si
podremos decidir sobre la terminación del sistema de reescritura. En
la segunda sección definiremos los órdenes de reducción, obteniendo
una propiedad con ellos para verificar la terminación. En el resto
daremos diferentes maneras de definir los órdenes de reducción.

\section{El problema de decisión}

\subsection{La indecibilidad para el caso general}

En esta sección vamos a introducir conceptos generales de las Ciencias
de la Computación, aplicados a los sistemas de reescritura. Uno de los
problemas generales de las Ciencias de la Computación es afirmar si un
problema es decidible. Usaremos algunos de sus resultados sobre
máquinas de Turing. Supondremos, sin perdida de generalidad, que el
modelo es no determinista de una banda finita en ambas direcciones.

\begin{defi} 
  Una máquina de Turing no determinista $M$ viene dada por,
  \begin{enumerate}
  \item Un alfabeto $\Gamma := \{s_0, \dots, s_n \}$ de símbolos,
    donde $s_0$ es considerado el espacio en blanco.
  \item Un conjunto finito $Q = \{q_0, \dots, q_p \}$ de estados.
  \item Una relación de transición
    $\Delta \subseteq Q \times \Gamma \times Q \times \Gamma \times \{
    l,r \} $.
  \end{enumerate}
\end{defi}
% Dibujo de una maquina de turing 

Una maquina de Turing se puede interpretar de la siguiente
manera. Imaginemos que tenemos una tira de papel dividida en infinitos
cuadrados en ambas direcciones, y en cada cuadrado se encuentra un símbolo del
alfabeto. Además de la tira de papel, se encuentra un marcador en uno
de los cuadrados. A partir de la relación de transición que escojamos,
el marcador se moverá de un lado a otro y cambiará los símbolos del
papel. 

Un ejemplo de relación de transición es $\{ q_1, s_1, q_2, s_2, l
\}$. El marcador empieza por el cuadrado inicial de la tira de
papel. Si el marcador posee el estado $q_1$ y en ese cuadrado de
papel se encuentra $s_1$, procedemos a aplicar el algoritmo. Cambiamos
el estado del marcador por $q_2$ y el símbolo del trozo de papel por
$s_1$. A continuación movemos el marcador a la izquierda (si fuera
$l$) o derecha (si fuera $r$).

%EJEMPLO

Llamaremos a un estado de la computación de la máquina,
configuración. En la configuración de una máquina se incluye, la tira
de papel, la posición y el estado del marcador. Podemos expresar
$K \vdash_M K'$ si la configuración de $K'$ se puede obtener mediante
$K$ con una computación de la máquina de Turing $M$.

De aquí surge el problema de la parada. Nos interesa saber si dado una
maquina de Turing $M$ y una configuración de la máquina $K$,
termina. Es decir, que no ocurre $K \vdash_M K_1 \vdash_M K_2 \dots$.

Volviendo al tema de la reescritura, en nuestro es caso el problema es
un poco más difícil. Como podemos aplicar toda regla en cada momento,
no tenemos una configuración $K$ como en el problema de la
parada. Pero si planteamos el problema suponiendo que no partimos de
una configuración inicial $K$, la pregunta pasa a ser; ¿qué
configuraciones $K$ hace que el problema termine o no? A esto se le
denota por el problema de la parada uniforme. 

Para extrapolar lo que ya sabemos sobre máquinas de Turing a la
reescritura, vamos a codificar las configuraciones como términos con
una signatura $\Sigma_M$.

\begin{defi}
Sea $M$ una máquina de Turing, definimos 
\[
  \Sigma_M := \{ \overrightarrow{s_0}, \dots, \overrightarrow{s_n},
  \overleftarrow{s_0}, \dots, \overleftarrow{s_n} \} \cup \{q_0,
  \dots, q_p \} \cup \{\overrightarrow{l}, r \}
\]
donde cada función es de aridad 1.
\end{defi}

A continuación, enunciamos una definición que relaciona términos con
máquinas de Turing.

\begin{defi}
  Sea $x_0$ una variable fija. Un término de configuración en
  $\Sigma_M$ es cualquier término de la forma,
\[
  \overrightarrow{l}(\overrightarrow{s_{i_k}}(\dots
  \overrightarrow{s_{i_1}}(q(\overleftarrow{s_{j_1}}(\dots
  \overrightarrow{s_{j_h}}(\overleftarrow{r}(x_0))\dots)))\dots))
\]
donde $k,h \geq 0$, $\{i_1, \dots, i_k, j_1, \dots, j_h \} \subseteq \{0, \dots, n \}$, y $q \in Q$.
\end{defi}

Para entender mejor esta definición, debemos hacer las siguientes
consideraciones. El marcador se encuentra en $s_{j_1}$ con el estado
$q$. Los elementos que se encuentran a la derecha del marcador son
$s_{j_2}, \dots s_{j_h}$, y los elementos que se encuentran a su
derecha, son $s_{i_1}, \dots, s_{i_k}$. 

Como la tira es infinita, $\overrightarrow{l}, \overleftarrow{r}$ nos
ayudan a controlar los espacios vacíos. Si durante la computación de
la máquina de Turing hiciese falta un espacio en blanco, estas
funciones se lo proporcionan.

%Insertar dibujo con la transformación

Acabamos de transformar una máquina de Turing en un término. En la
siguiente definición, adaptaremos un sistema de reescritura para una
máquina de Turing.

\begin{defi}
  Un sistema de reescritura $R_M$ consiste en las siguientes reglas de
  reescritura,
\begin{itemize}
\item Para cada transición $(q, s_i, q', s_j, r) \in \Delta$, $R_M$
  contiene la regla,
\[
  q(\overleftarrow{s_i}(x)) \rightarrow \overrightarrow{s_j}(q'(x))
\]
Si $i=0$, entonces $R_M$ contiene la siguiente regla,
\[
  q(\overleftarrow{r}(x)) \rightarrow
  \overrightarrow{s_j}(q'(\overleftarrow{r}(x)))
\]
\item Para cada transición $(q,s_i,q',s_j,l) \in \Delta$, $R_M$
  contiene la regla,
\[
  \overrightarrow{l}(q(\overleftarrow{s_i}(x))) \rightarrow
  \overrightarrow{l}(q'(\overleftarrow{s_0}(\overleftarrow{s_j}(x))))
\]
y para cada $s_k \in \Gamma$ la regla,
\[
  \overrightarrow{s_k}(q(\overleftarrow{s_i}(x))) \rightarrow
  q'(\overleftarrow{s_k}(\overleftarrow{s_j}(x)))
\]
Si $i=0$, entonces $R_M$ contiene una regla adicional,
\[
  \overrightarrow{l}(q(\overleftarrow{r}(x))) \rightarrow
  \overrightarrow{l}(q'(\overleftarrow{s_0}(\overleftarrow{s_j}(\overleftarrow{r}(x)))))
\]
y para cada $s_k \in \Gamma$ la regla
\[
  \overrightarrow{s_k}(q(\overleftarrow{r}(x))) \rightarrow
  q'(\overleftarrow{s_0}(\overleftarrow{s_j}(\overleftarrow{r}(x)))))
\]
\end{itemize}
\end{defi}

Con estas reglas podemos llegar a la conclusión de que para todo par
de términos de configuración $t, t'$, que verifiquen
$t \rightarrow_{R_M} t'$, implica $K_t \vdash K_{t'}$. E igual ocurre
en sentido contrario, si tenemos dos configuraciones $K, K'$ y un
término de configuración $t$, si $K \vdash K'$ y $K \equiv K_t$, esto
implica que existe un término de configuración $t'$ tal que
$K' \equiv K_{t'}$ y $t \rightarrow_{R_M} t'$. \label{prop:4}

De aquí podemos razonar una propiedad. Como el problema de la parada
de las máquinas de Turing es indecidible, y acabamos de probar que una
máquina de Turing es equivalente a un sistema de reescritura (con las
reglas que hemos pedido), aseguramos que dado un sistema de
reescritura finito $R$ y un término $t$, el problema de ver si todas
las reducciones son terminantes empezando por $t$ es indecidible.

Sin embargo, el problema que acabamos de resolver no es el original
que queríamos. El problema de la terminación pide que todas las
reducciones desde todos los posibles términos sean terminantes. Puede
darse el caso que tengamos una configuración de términos que termine,
pero una reducción no terminante que empiece por un término que no
este en la configuración inicial. Por tanto el problema no esta
resuelto todavía. El siguiente lema asegura que este caso no puede
ocurrir.

\begin{lema} \label{lema:4.1}
  Sea  $t$  un  término  de   $\Sigma_M$.   Si  existe  una  reducción
  $t \rightarrow_{R_M}  t_1 \rightarrow_{R_M} t_2  \rightarrow \dots$,
  entonces existe  un término  de configuración  $t'$ y  una reducción
  infinita $R_M$ empezando por $t'$.
\end{lema}

\begin{demo}
  Vamos a considerar
  $\Sigma_M = \overrightarrow{\Gamma} \cup \overleftarrow{\Gamma} \cup
  Q \{ \overleftarrow{r} , \overrightarrow{l} \} $, donde
  $\overrightarrow{\Gamma} := \{\overrightarrow{s_1}, \dots,
  \overrightarrow{s_n} \}$ y
  $\overleftarrow{\Gamma} := \{ \overleftarrow{s_1}, \dots,
  \overleftarrow{s_n} \}$. Cualquier elemento $w$ de $\Sigma_M$ puede
  ser escrito como composición de funciones tal que,
  $w = u_1 v_1 u_2 v_2 \dots u_q v_q u_{q+1}$, ya que $u_i, v_j$ son
  funciones de aridad 1.

  Como todas las reglas de reescritura de $R_M$ contienen un símbolo
  de $Q$, cualquier reducción aplicada a $w$, se hace dentro de las
  funciones $v$.

  Es decir, que si tomamos $w(x) \rightarrow_{R_M} w'(x)$, entonces
  existe un índice $j$, $1 \leq j \leq q$, y una palabra
  $v'_j \in \overrightarrow{\Gamma}^{*} Q \overleftarrow{\Gamma}^{*}$
  tal que
  $w'= u_1 v_1 u_2 v_2 \dots u_j v'_j u_{j+1} \dots u_q v_1 u_{q+1}$,
  y que
  $\overrightarrow{l} v_j \overleftarrow{r} (x_0) \rightarrow_{R_M}
  \overrightarrow{l} v'_j \overleftarrow{r}$
  
  Como $q$ es finito, esto implica que la reducción infinita que
  empieza por $w(x)$ produce una reducción infinita empezando por
  $\overrightarrow{l} v_j \overleftarrow{r} (x_0)$. Como
  $\overrightarrow{l} v_j \overleftarrow{r} (x_0)$ es un término de
  configuración, queda demostrado el lema.
\end{demo}

Finalmente, por el lema \ref{lema:4.1}, obtenemos el objetivo principal
de esta sección, la indecibilidad para el problema de terminación de
los sistemas de reescritura.


\subsection{Sistemas de reescritura básicos hacia la derecha}

En esta subsección vamos a analizar un SRT particular $R$, donde sus
reglas hacia la derecha son términos básicos. A estos SRT les
llamaremos básicos hacia la derecha, y probaremos que son terminantes.

\begin{lema}
  Sea $R$ un SRT finito básico hacia la derecha. Entonces son
  equivalentes,
  \begin{enumerate}
  \item $R$ no termina.
  \item Existe una regla $l \rightarrow r \in R$ y un término $t$ de
    manera que $r \xrightarrow{+}_R t$ y $t$ contiene a $r$ como
    subtérmino.
  \end{enumerate}
\end{lema}

\begin{demo}
  $(2 \Rightarrow 1)$ es cierto pues obtenemos una reducción finita;
  $r\xrightarrow{+}_R t = t[r]_p \xrightarrow{+}_R t[t]_p =
  t[t[r]_p]_p \xrightarrow{+}_R \dots$, donde $p$ es la posición $t|_p = r$.
  
  $(1 \Rightarrow 2)$ se demuestra mediante inducción por el cardinal
  de $R$.

  Suponemos que $|R| >0>$ y consideramos una reducción infinita
  $t_1 \rightarrow t_2 \rightarrow t_3 \rightarrow \dots$. Tenemos que
  probar que esta cadena no termina. Sin perdida de generalidad al
  menos una reducción ocurre en la posición $\epsilon$. Esto significa
  que existe un índice $i$, una regla $l \rightarrow r \in R$, y una
  sustitución $\sigma$, tal que $t_i = \sigma(l)$ y
  $t_{i+1} = \sigma(r) = r$. Luego hay una reducción infinita
  $r \rightarrow_R t_{i+2} \rightarrow_R t_{i+3} \rightarrow_R \dots$
  que empieza por $r$.

  Si la regla $l \rightarrow r$ no es usada en la reducción, entonces
  aplicamos la inducción a $R - \{l \rightarrow r\}$. En el caso en
que sí sea usada, eso significa que existe $j$ tal que $r$ ocurre en
$t_{i+j}$ y por tanto, obtenemos la segunda proposición.
\end{demo}

Luego si $R$ termina entonces $\rightarrow_R$ es globalmente finita,
ya que es una ramificación finita por el corolario \ref{globfin} y el
lema \ref{ramfin}. Por tanto todas las reducciones terminarían y, al
ser finitas, lo hacen en un numero finito de pasos. Con esta idea,
obtenemos el resultado clave de esta subsección.

\begin{teor}
  Para un SRT finito básico hacia la derecha, el problema de la
  terminación es decidible.
\end{teor}

\section{Órdenes de reducción}

En esta sección relacionaremos la definición de orden bien fundado con
el problema de la terminación.

En la sección previa, hemos probado que el problema de la terminación
es indecidible. Sin embargo podemos dar varios métodos para resolver
el problema. Estos métodos no son totalmente automatizados.

La idea básica para probar la terminación, se hace mediante un orden
bien fundado. Supongamos que $R$ es un sistema de reescritura finito,
$>$ es un orden estricto bien fundado en $T(\Sigma, X)$. Relacionando
los sistemas de reescritura con los órdenes, si $R$ es terminante,
diremos que $s \rightarrow_R t$ implica $s > t$. En vez de decidir $s > t$,
tan solo tendremos que comprobar las reglas de $R$. Para considerar
este nuevo acercamiento al problema necesitamos pedir ciertas
propiedades a $>$.

\begin{defi}
  Sea $\Sigma$ una signatura y $V$ un conjunto numerable de
  variables. Un orden estricto $>$ es un orden de reescritura syss
  \begin{enumerate}
  \item Es compatible con las $Sigma$-operaciones: para todo
    $s_1,s_2 \in T(\Sigma, V)$, todo $n \geq 0$ y
    $f \in \Sigma^{(n)}$, $s_1 > s_2$ implica
    \[ f(t_1, \dots, t_{n-1}, s_1, t_{i+1}, \dots, t_n) > f(t_1,
      \dots, t_{n-1}, s_2, t_{i+1}, \dots, t_n) \]
    $\forall i, 1 \leq i \leq n$ y
    $\forall t_1, \dots, t_{i-1}, t_{i+1}, \dots, t_n \in T(\Sigma,
    V)$. %Hay una condición equivalente
  \item Es cerrado bajo sustituciones:
    $\forall \sigma \in \Sub(T(\Sigma,V))$, si $s_1 > s_2$ entonces
    $\sigma(s_1) > \sigma(s_2)$.
  \end{enumerate}
  Un orden de reducción es un orden de reescritura bien fundado.
\end{defi}

%Ejemplo

El siguiente teorema es de especial importancia para los órdenes de
reducción.

\begin{teor} 
  Un sistema de reescritura $R$ termina syss existe un orden de
  reducción $>$ que satisface $l>r$ para toda regla
  $l \rightarrow r \in R$.
\end{teor}

\begin{demo}
  $(\Rightarrow)$ Suponiendo que $R$ termina, entonces
  $\xrightarrow{+}_R$ es un orden de reducción ya que satisface
  $l \xrightarrow{x}_R r$ para todo $l \rightarrow r \in R$.

  $(\Leftarrow)$ Como $>$ es un orden de reescritura, para $l > r$ se
  verifica $t [ \sigma(l) ]_p > t [ \sigma(r) ]_p$ para todo
  término $t$, sustitución $\sigma$ y posición $p$. Como $>$ esta bien
  fundada, y para toda regla de $R$, $s_1 \rightarrow_R s_2$ implica
  $s_1 > s_2$, entonces no puede ocurrir una reducción infinita
  $s_1 \rightarrow_R s_2 \rightarrow_R s_3 \dots$.
\end{demo}

%\begin{defi}
%  Un orden de reducción $>$ en $T(\Sigma, V)$ es total en términos
%  básicos, si restringido a $T(\Sigma, \emptyset)$ es un orden linear
%  estricto.
%\end{defi}

A partir de aquí podemos dar varios métodos para dar una repuesta al
problema. En las secciones posteriores nos centraremos en los ordenes
de simplificación, orden de camino lexicográfico y orden de camino
recursivo, y daremos una implementación de ambos.

\section{Órdenes de simplificación}

En esta sección definiremos que son los órdenes de simplificación y
daremos dos órdenes construidos a partir de estos a modo de ejemplo en
las siguientes subsecciones. Empezamos por la definición,

\begin{defi}
  Sea $>$ un orden estricto en $T(\Sigma, V)$ es un orden de
  simplificación, syss es un orden de reescritura que para todo
  término $t \in T(\Sigma, V)$ y todas las posiciones
  $p \in \Pos(t) - \{ \epsilon \}$ entonces $t > t|_p$.
\end{defi}

La definición es equivalente a pedir que, para todo $n \leq 1$, todos
los símbolos de funciones $f \in \Sigma^{(n)}$, todas las variables
$x_1, \dots, x_n \in V$, tenemos que
$f(x_1, \dots, x_i, \dots, x_n) > x_i$. Usando esta nueva definición e
introduciendo algunos conceptos de álgebra de homomormismos, se puede
demostrar que los órdenes de simplificación son equivalentes a los
ordenes de reducción para $\Sigma$ finito.

\subsection{Órdenes de caminos lexicográficos}

La idea de los caminos recursivos reside en comparar las raíces de los
términos, y en comparar recursivamente sus subtérminos. Estos
subtérminos se pueden comparar mediante multiconjuntos (orden de
caminos de multiconjuntos), tuplas (orden de caminos lexicográficos),
o una mezcla de ambos (orden de caminos recursivos). Nos interesaremos
por estos dos últimos.

\begin{defi}
  Sea $\Sigma$ una signatura finita y $>$ un orden estricto de
  $\Sigma$. El orden de caminos lexicográficos $>_{ocl}$, se define
  como, $s>_{ocl}t$ syss ocurre alguno de los siguientes casos,
  \begin{itemize}
  \item (OCL1) $t \in \Var(s)$ y $s \not = t$ ó
  \item (OCL2) $s = f(s_1, \dots, s_m)$, $t = g(t_1, \dots, t_n)$, y
    \begin{itemize}
    \item (OCL2a) existe $i$, $1 \leq i \leq m$, con
      $s_i \geq_{lpo} t$ ó
    \item (OCL2b) $f > g$ y $s >_{ocl} t_j$ para todo $j$ en
      $1 \leq j \leq n$ ó
    \item (OCL2c) $f = g$, $s >_{lpo} t_j$ para todo $j$ en
      $1 \leq j \leq n$, y exista $i$, $1 \leq i \leq m$ tal que
      $s_1 = t_1, \dots, s_{i-1} = t_{i-1}$ y $s_i >_{ocl} t_i$
    \end{itemize}
  \end{itemize}
\end{defi}

La definición es recursiva y esta bien definida.

% INTRODUCIR EJEMPLO

\begin{teor}
  Para cualquier orden estricto $>$ en $\Sigma$, el orden de camino
  lexicográfico inducido es un orden de simplificación en
  $T(\Sigma, V)$.
\end{teor}

% FALTA DEMOSTRACIÓN

A continuación implementaremos el orden de caminos lexicográfico en
Haskell. Para realizar los ejemplos de esta función, primero
implementaremos \texttt{ordenPorLista xs a b}, que es el resultado de
comparar $a$ y $b$, tal que $a > b$ syss $a$ aparece antes en $x$s que
$b$. Por ejemplo,

\begin{sesion}
ghci> ordenPorLista ["a","b","c"] "a" "c"
GT
ghci> ordenPorLista ["a","b","c"] "c" "b"
LT
ghci> ordenPorLista ["a","b","c"] "b" "b"
EQ
\end{sesion}

Su código es,

\begin{codigo}
ordenPorLista :: Ord a => [a] -> a  -> a -> Ordering
ordenPorLista [] _ _ =
  error("Ninguno de los dos elementos se encuentran
         en la lista")
ordenPorLista (x:xs) a b
  | a == x = if a == b
             then EQ
             else GT
  | b == x = LT
  | otherwise = ordenPorLista xs a b
\end{codigo}

Definiremos \texttt{ordCamLex ord s t}, que es el resultado de
comparar $s$ y $t$ términos, con el orden de caminos lexicográfico
inducido por \texttt{ord}. Por ejemplo

\begin{sesion}
ghci> let ord = ordenPorLista ["i","f","e"]
ghci> ordenCamLex ord (T "f" [V ("x",1), T "e" []]) (V ("x",1))
GT
ghci> ordenCamLex ord (T "i" [T "e" []]) (T "e" [])
GT
ghci> ordenCamLex ord (T "i" [T "f" [V("x",1),V("y",1)]])
                      (T "f" [T "i" [V("y",1)], T "i" [V("x",1)]])
GT
ghci> ordenCamLex ord (T "f" [V("y",1),V("z",1)])
                      (T "f" [T "f" [V("x",1),V("y",1)], V("z",1)])
LT
\end{sesion}

Su código es,

\begin{codigo}
 ordenCamLex:: ([Char] -> [Char] -> Ordering)
                -> Termino -> Termino -> Ordering
ordenCamLex _ s (V x)
  | s == (V x) = EQ
  | ocurre x s = GT --OCL1
  | otherwise = LT
ordenCamLex _ (V _) (T _ _) = LT
ordenCamLex ord s@(T f ss) t@(T g ts) --OCL2
  | all (\x -> ordenCamLex ord x t == LT) ss
    = case ord f g of
      GT -> if all (\x -> ordenCamLex ord s x == GT) ts
            then GT --OCL2b
            else LT
      EQ -> if all (\x -> ordenCamLex ord s x == GT) ts
            then ordLex (ordenCamLex ord) ss ts --OCL2c
            else LT
      LT -> LT  
  | otherwise = GT --OCL1a
\end{codigo}

\subsection{Órdenes de caminos recursivos}


%%%%%%%%%%%%%%%%%%%%%%%%%%%%%%%%%%%%%%%%%%%%%%%%%%%
%%Apéndices
%%%%%%%%%%%%%%%%%%%%%%%%%%%%%%%%%%%%%%%%%%%%%%%%%%%

\clearpage
\addappheadtotoc
\appendix


%%% Local Variables:
%%% mode: latex
%%% TeX-master: "SRT_en_Haskell"
%%% End:

\chapter{Confluencia}

En este capítulo estudiaremos el problema de determinar si un sistema
de reescritura es confluente. En la primera sección vamos a demostrar
que este problema es indecidible para el caso general. Sin embargo en
las secciones posteriores, estudiaremos que ocurre para un sistema
terminante. Para este caso el problema será decidible. Por último
veremos que podemos decir para los sistemas que no terminan.

\section{Estudio sobre el problema de decisión}
  
En esta sección veremos la indecibilidad para comprobar si un sistema
es confluente mediante el siguiente resultado,

\begin{teor}
  El problema de decidir si un sistema de reescritura finito $R$ es
  confluente, es indecidible.
\end{teor}

\begin{demo}
  El objetivo de esta demostración es reducir el problema de las
  palabras básicas para $E$ (que sabemos que es indecidible), a un
  sistema de reescritura de términos.

  Sea un conjunto de identidades $E$ tal que $\Var(l) = \Var(r)$ para
  todo $l \approx r \in E$. Sea $R := E \cup E^{-1}$, entonces al
  tener $\rightarrow_R = \leftrightarrow_E$ es confluente. Además $R$
  es un sistema de reescritura (por $\Var(l) = \Var(r)$).

  Dados dos términos básicos $s$ y $t$, y una constante $a$, vamos a
  probar que $R_{st} := R \cup \{ a \rightarrow s, a \rightarrow t \}$
  es confluente syss $s \approx_E t$.

  \begin{itemize}
  \item $(\Rightarrow)$ Si $R_{st}$ es confluente, entonces ni $s$, ni
    $t$ poseen una constante $a$, luego $s \downarrow_{R_{st}} t$. Esto
    significa que las reglas $a \rightarrow s, a \rightarrow t$ no
    pueden ser usadas. Por tanto $\rightarrow_R = \leftrightarrow_E$ y
    $s \approx_E t$.

  \item $(\Rightarrow)$ Vamos a probar que para términos $u, v$, si
    $u \rightarrow_{R_{st}} v$ entonces $v^t \xrightarrow{*}_R u^t$,
    donde $u^t$ denota el resultado de sustituir las constantes $a$ en
    $u$ por $t$.

    Supongamos que $u \rightarrow_{R{_st}} v$. Distinguimos que reglas
    se usan.
    \begin{itemize}
    \item Si usamos $u \rightarrow_R v$, reemplazamos $a$ por $t$ para
      conseguir $u^t \rightarrow_R v^t$ y por tanto
      $v^t \rightarrow_R u^t$ al ser $R$ simétrico.
    \item Si usamos $a \rightarrow s$, entonces $u|_p = a$ y
      $v = u[s]_p$ para alguna posición $p$. Como $s \approx_E t$
      entonces $s \xrightarrow{*}_R t$. Obtenemos
      $v \xrightarrow{*}_R u[t]_p$ que conlleva a
      $v^t \xrightarrow{*}_R (u[t]_p)^t$. Pero $(u[t]_p)^t =
      u^t[t^t]_p = u^t[t]_p = u^t$.
    \item Si usamos $a \rightarrow t$ en la posición $p$, entonces
      $v^t = (u[t]_p)^t = u^t \xrightarrow{*}_R u^t$.
    \end{itemize}
    Por tanto, si $u \xrightarrow{*}_{R_{st}} u_i$, para $i= 1,2$,
    entonces $u_i \xrightarrow{*}_{R_{st}} u_i^t \xrightarrow{*}_R u^t$
    y obtenemos $u_1 \downarrow_{R_{st}}$
  \end{itemize}
\end{demo}

\section{Pares críticos}

En esta sección estudiaremos la decibilidad para sistemas de
reescritura finitos que sean localmente confluentes.

La necesidad de los pares críticos surge de la importancia a la hora
de aplicar las reglas de un sistema de reescritura. Para entender
mejor este concepto, plantearemos el siguiente ejemplo. Sea $s$ un
término, $R := \{ l_1 \rightarrow r_1, l_2 \rightarrow r_2 \}$ donde
$l_1, l_2$ son subtérminos de $s$. Según apliquemos la primera regla o
la segunda, obtendremos un nuevo término $t$.

\begin{figure}[h]
  \centering
  \begin{tikzpicture}[->,>=stealth',level/.style={sibling distance =
      5cm/#1, level distance = 1.5cm}]

\node{$s$}
    child {
      node[align=center]{$t_1$}
        edge from parent node[above left] {$l_1 \rightarrow r_1$}}
    child {
      node{$t_2$}
      edge from parent node[above right] {$l_2 \rightarrow r_2$}
    };
\end{tikzpicture}
\end{figure}

Sea $p_i$ las posiciones y $\sigma_i$ las sustituciones tal que
$s|_{p_i} = \sigma_i l_i$ y $t_i = s[\sigma_i r_i ]_{p_i}, i
=1,2$. Tenemos que estudiar como de relacionados están $p_1$ y $p_2$.

\begin{itemize}

\item \textbf{Caso 1}, $p_1$ y $p_2$ están en árboles separados. En este caso,
  se da la convergencia local. Por ejemplo, supongamos que partimos de
  la figura \ref{fig:caso1a}.

  \begin{figure}[h]
    \centering
    \begin{tikzpicture}
\draw (0,0) node[anchor=north]{}
  -- (10,0) node[anchor=north]{}
  -- (5,5) node[below]{}
  -- cycle
   (2,0) node[anchor=north]{}
  -- (4,0) node[anchor=north]{}
  -- (3,1) node[below]{}
  -- cycle
  
     (8,0) node[anchor=north]{}
  -- (6,0) node[anchor=north]{}
  -- (7,1) node[below]{}
  -- cycle
  ;
\node[below] at (3,0.8) {\LARGE $\sigma_1 l_1$};
\node[below] at (7,0.8) {\LARGE $\sigma_2 l_2$};
\node[below] at (5,4.7) {\huge $s$};
\end{tikzpicture}
\caption{Término del caso 1 sin aplicar ninguna regla}
\label{fig:caso1a}
\end{figure}

Si aplicamos $l_1 \rightarrow r_1$ o $l_2 \rightarrow r_2$ obtenemos
las figuras \ref{fig:caso1b} y \ref{fig:caso1c} respectivamente.

\begin{figure}[h]
    \centering
\begin{tikzpicture}
\draw (0,0) node[anchor=north]{}
  -- (10,0) node[anchor=north]{}
  -- (5,5) node[below]{}
  -- cycle
   (2,0) node[anchor=north]{}
  -- (4,0) node[anchor=north]{}
  -- (3,1) node[below]{}
  -- cycle
  
     (8,0) node[anchor=north]{}
  -- (6,0) node[anchor=north]{}
  -- (7,1) node[below]{}
  -- cycle
  ;
\node[below] at (3,0.8) {\LARGE $\sigma_1 l_1$};
\node[below] at (7,0.8) {\LARGE $\sigma_2 l_2$};
\node[below] at (5,4.7) {\huge $s$};

\node(a) at (2,0) {};
\node (b) at (4,0) {};
\node (c) at (3,1) {};
\fill (a.center) -- (b.center) -- (c.center);

\end{tikzpicture}
\caption{Término del caso 1 tras aplicar $l_1 \rightarrow r_1$}
\label{fig:caso1b}
\end{figure}

\begin{figure}[h]
\centering
\begin{tikzpicture}
\draw (0,0) node[anchor=north]{}
  -- (10,0) node[anchor=north]{}
  -- (5,5) node[below]{}
  -- cycle
   (2,0) node[anchor=north]{}
  -- (4,0) node[anchor=north]{}
  -- (3,1) node[below]{}
  -- cycle
  
     (8,0) node[anchor=north]{}
  -- (6,0) node[anchor=north]{}
  -- (7,1) node[below]{}
  -- cycle
  ;
\node[below] at (3,0.8) {\LARGE $\sigma_1 l_1$};
\node[below] at (7,0.8) {\LARGE $\sigma_2 l_2$};
\node[below] at (5,4.7) {\huge $s$};

\node(a) at (8,0) {};
\node (b) at (6,0) {};
\node (c) at (7,1) {};
\fill (a.center) -- (b.center) -- (c.center);

\end{tikzpicture}
\caption{Término del caso 1 tras aplicar $l_2 \rightarrow r_2$}
\label{fig:caso1c}
\end{figure}

Y si usamos la regla que no se ha aplicado antes, para ambos términos
llegamos a la figura \ref{fig:caso1d}.

\begin{figure}[h]
\centering
\begin{tikzpicture}
\draw (0,0) node[anchor=north]{}
  -- (10,0) node[anchor=north]{}
  -- (5,5) node[below]{}
  -- cycle
   (2,0) node[anchor=north]{}
  -- (4,0) node[anchor=north]{}
  -- (3,1) node[below]{}
  -- cycle
  
     (8,0) node[anchor=north]{}
  -- (6,0) node[anchor=north]{}
  -- (7,1) node[below]{}
  -- cycle
  ;
\node[below] at (3,0.8) {\LARGE $\sigma_1 l_1$};
\node[below] at (7,0.8) {\LARGE $\sigma_2 l_2$};
\node[below] at (5,4.7) {\huge $s$};

\node(a) at (8,0) {};
\node (b) at (6,0) {};
\node (c) at (7,1) {};
\node(a2) at (2,0) {};
\node (b2) at (4,0) {};
\node (c2) at (3,1) {};
\fill (a.center) -- (b.center) -- (c.center);
\fill (a2.center) -- (b2.center) -- (c2.center);

\end{tikzpicture}
\caption{Término en el caso 1, donde no se produce solapamiento}
\label{fig:caso1d}
\end{figure}

  
\item \textbf{Caso 2}. Si $p_1$ es un prefijo de $p_2$, tal y como se
  muestra en la figura \ref{fig:caso2}

\begin{figure}[h]
  \centering
\begin{tikzpicture}
\draw[dotted](5,3)--(10,3);
\draw[dotted](5,5)--(10,5);
\draw[dotted](5,1)--(10,1);

\draw [decorate,decoration={brace,amplitude=10pt,mirror,raise=4pt},yshift=0pt]
(10,1) -- (10,3) node [black,midway,xshift=0.8cm] {};

\draw [decorate,decoration={brace,amplitude=10pt,mirror,raise=4pt},yshift=0pt]
(10,3) -- (10,5) node [black,midway,xshift=0.8cm] {};

\draw [decorate,decoration={brace,amplitude=10pt,mirror,raise=4pt},yshift=0pt]
(11,1) -- (11,5) node [black,midway,xshift=0.8cm] {};


\draw (0,0) node[anchor=north]{}
  -- (10,0) node[anchor=north]{}
  -- (5,5) node[below]{}
  -- cycle
   (2,0) node[anchor=north]{}
  -- (8,0) node[anchor=north]{}
  -- (5,3) node[below]{}
  -- cycle
  
     (4,0) node[anchor=north]{}
  -- (6,0) node[anchor=north]{}
  -- (5,1) node[below]{}
  -- cycle
  ;

\node[below] at (5,4.7) {\huge $s$};
\node[below] at (5,2.7) {\LARGE $\sigma_1 l_1$};
\node[below] at (5,0.8) {\LARGE $\sigma_2 l_2$};
\node[below] at (10.9,4.3) {\LARGE $p_1$};
\node[below] at (10.8,2.3) {\LARGE $p$};
\node[below] at (11.9,3.3) {\LARGE $p_2$};
\end{tikzpicture}
\caption{Situación del caso 2.}
\label{fig:caso2}
\end{figure}


A partir del dibujo, podemos distinguir dos casos según como de
separadas estén los términos $\sigma_1 l_1$ y $\sigma_2 l_2$.

  \begin{itemize}

  \item \textbf{Caso 2.1}. La reducción $\sigma_2 l_2$ no se solapa
    con $l_1$ pero está contenido en $\sigma_1$ (figura
    \ref{caso2a}). A esta situación la llamaremos solapamiento no
    crítico, pues en esta situación ocurre una confluencia local.

\begin{figure}[h]
  \centering
\usetikzlibrary{decorations.pathmorphing}
  \begin{tikzpicture}
\draw[dotted](5,3)--(10,3);
\draw[dotted](5,5)--(10,5);
\draw[dotted](5,2)--(10,2);

\draw(3,3)--(7,3);

\draw [decorate,decoration={brace,amplitude=10pt,mirror,raise=4pt},yshift=0pt]
(10,2) -- (10,3) node [black,midway,xshift=0.8cm] {};

\draw [decorate,decoration={brace,amplitude=10pt,mirror,raise=4pt},yshift=0pt]
(10,3) -- (10,5) node [black,midway,xshift=0.8cm] {};

\draw [decorate,decoration={brace,amplitude=10pt,mirror,raise=4pt},yshift=0pt]
(11,2) -- (11,5) node [black,midway,xshift=0.8cm] {};

\draw (0,0) node[anchor=north]{}
  -- (10,0) node[anchor=north]{}
  -- (5,5) node[below]{}
  -- cycle
  
     (3,0) node[anchor=north]{}
  -- (7,0) node[anchor=north]{}
  -- (5,2) node[below]{}
  -- cycle
  ;
  
\node[below] at (5,4.3) {\huge $l_1$};
\node[below] at (5,2.8) {\LARGE $\sigma_1$};
\node[below] at (5,1.2) {\LARGE $\sigma_2 l_2$};
\node[below] at (10.9,4.3) {\LARGE $q_1$};
\node[below] at (10.9,2.8) {\LARGE $q_2$};
\node[below] at (11.9,3.3) {\LARGE $p$};
\end{tikzpicture}
\caption{Caso 2.1.}
\label{fig:caso2a}
\end{figure}
    
\item \textbf{Caso 2.2}. En este caso $l_1$ y $l_2$ se solapan (figura
  \ref{caso2b}). A este caso lo llamaremos solapamiento crítico, y es
  la peor situación que podemos tener.

 \begin{figure}
      \centering
\begin{tikzpicture}


\draw[dotted](5,3)--(10,3);
\draw[dotted](5,5)--(10,5);


\draw(4,2)--(2,2);
\draw(6,2)--(8,2);
\draw(3.5,1.5)--(6.5,1.5);


\draw [decorate,decoration={brace,amplitude=10pt,mirror,raise=4pt},yshift=0pt]
(10,3) -- (10,5) node [black,midway,xshift=0.8cm] {};


\draw (0,0) node[anchor=north]{}
  -- (10,0) node[anchor=north]{}
  -- (5,5) node[below]{}
  -- cycle
  
     (2,0) node[anchor=north]{}
  -- (8,0) node[anchor=north]{}
  -- (5,3) node[below]{}
  -- cycle
  ;
  
\node[below] at (5,4.3) {\huge $l_1$};
\node[below] at (5,2.5) {\LARGE $l_2$};
\node[below] at (5,1) {\LARGE $\sigma_2$};
\node[below] at (2.25,1.5) {\LARGE $\sigma_1$};
\node[below] at (7.75,1.5) {\LARGE $\sigma_1$};
\node[below] at (10.9,4.3) {\LARGE $p$};

\end{tikzpicture}

\caption{Caso 2.2.}
\label{caso2b}
    \end{figure}    
  \end{itemize}


\end{itemize}


\begin{defi}
  Sea $l_i \rightarrow r_i$, $i = 1,2$ dos reglas cuyas variables han
  sido renombradas tal que
  $\Var(l_1,r_1) \cap \Var(l_2,r_2) = \emptyset$. Sea
  $p \in \Pos(l_1)$ tal que $l_1|_p$ no es una variable, y $\theta$ un
  umg de $l_1|_p =^? l_2$. Esto determinará el par crítico $\langle
  \theta r_1, (\theta l_1) [\theta r_2 ]_p \rangle$.
  % DIBUJO DEL ARBOL
  Si dos reglas dan lugar a un par crítico, diremos que se solapan.
\end{defi}

% EJEMPLO

A continuación enunciamos el Teorema de los Pares Críticos. Este
resultado así como la idea de la demostración, también se usará en la
siguiente sección.
\begin{teor}
  Un sistema de reescritura es localmente confluente syss todos sus
  pares críticos se vuelven a unir.
\end{teor}

\begin{demo}
  ($\Leftarrow$) Sabemos que si $s \rightarrow_R t$ para $i = 1,2$,
  entonces es localmente confluente, ó $t_i = s[u_i]_p$, donde
  $\langle u_1, u_2 \rangle$ proviene de algún par crítico
  $\langle v_1, v_2 \rangle$, es decir $u_1 = \delta v_i$. Entonces
  $\xrightarrow{*}t$ para algún término, y por tanto
  $u_i \xrightarrow \delta t$ es también cierto. Esto implica
  $t_i \xrightarrow s[\delta t]_p$ para $i = 1,2$.

  ($\Rightarrow$) Como el sistema de reescritura es localmente
  confluente, esto significa que cuando lleguemos a un par crítico
  este volverá a converger a un término. Luego todos sus pares
  críticos se vuelven a unir.
\end{demo}

Algunas propiedades directas de este teorema,

\begin{coro}\label{parcritcor}
  Un sistema de términos terminante es confluente syss todos sus pares
  críticos se vuelven a unir
\end{coro}

Pero si además pedimos que sea un sistema finito, obtenemos el
siguiente resultado.

\begin{coro}
  La confluencia de un sistema de reescritura finito y terminante es
  decidible.
\end{coro}

\begin{demo}
  El algoritmo que seguiremos será el siguiente:

  Para cada par de reglas $l_1 \rightarrow r_1$ y
  $l_2 \rightarrow r_2$, y para cada $p \in \Pos(l_1)$, tal que
  $l_1|_p$ no es una variable, se genera un par crítico unificando las
  variables disjuntas de $l_1|_p$ y $l_2$.

  Para cada uno de esos pares críticos $\langle u_1, u_2 \rangle$
  reducimos $u_i$ a su forma normal $\tilde{u}_i$. Decimos que el
  sistema de reescritura es confluente syss
  $\tilde{u}_1 = \tilde{u}_2$ para todos sus pares críticos.

  Si se verifica $\tilde{u}_1 = \tilde{u}_2$, entonces por el
  corolario \ref{parcritcor}, el sistema es confluente.

  Si se da el caso de que existe algún par crítico no
  $\tilde{u}_1 \not = \tilde{u}_2$, entonces se da la situación
  $\tilde{u}_1 \xleftarrow{*} u_1 \leftarrow u \rightarrow u_2
  \xrightarrow{*} \tilde{u_2}$, que demostraría que no es confluente.
\end{demo}

\section{Implementación de los pares críticos}

Se implementarán una función para calcular los pares críticos. Antes
de poder definirla, se necesitarán varias funciones auxiliares. El
código se puede encontrar en el fichero \texttt{ParesCriticos.hs}.

Vamos a necesitar la librería de términos.
\begin{codigo}
import Terminos
\end{codigo}

\begin{itemize}
\item \texttt{(indiceMaximo t)} \index{\texttt{indiceMaximo}}es el mayor índice que aparece en
  \texttt{t}. Por ejemplo,
\begin{sesion}
>>> indiceMaximo (V ("x",3))
3
>> indiceMaximo (T "w" [V("x",4), V("z", 200), V("y", 1)])
200
\end{sesion}

Su código es,

\begin{codigo}
indiceMaximo :: Termino -> Indice
indiceMaximo (V (_,i)) = i
indiceMaximo (T _ ts) = maximum (map indiceMaximo ts)
\end{codigo}

\item \texttt{(renombraTermino n t)}\index{\texttt{renombraTermino}}
  es el termino resultante tras sumar a todos los índices de
  \texttt{t} el entero \texttt{n}. Por ejemplo,
\begin{sesion}
>>> renombraTermino 5 (V ("x",3)) 
V ("x",8)
>>> renombraTermino 3 (T "w" [V("x",4), V("z", 200), V("y", 1)])
T "w" [V ("x",7),V ("z",203),V ("y",4)]
\end{sesion}

Su código es,

\begin{codigo}
renombraTermino :: Int -> Termino -> Termino
renombraTermino n (V (x,i)) = V(x,i+n)
renombraTermino n (T f ts) = T f (map (renombraTermino n) ts)
\end{codigo}

\item \texttt{(parCritico c (l1,r1) (l2,r2)} \index{\texttt{parCritico}} calcula el par crítico de la
  regla $\texttt{l1} \rightarrow \texttt{r1}$ y
  $\texttt{l2} \rightarrow \texttt{r2}$. Por ejemplo,
\begin{sesion}
>>> let c a = a
>>> parCritico c (V ("x",1), T "f" [V ("y",1)])
                 (V ("x",2),T "g" [V("y",2)])
[(T "f" [V ("y",1)],T "g" [V ("y",2)])]
>>> parCritico c (T "f" [V ("x",1)],T "f" [V ("y",1)])
               (V ("x",1), T "g" [V("y",2)])
[]
\end{sesion}

Su código es,

\begin{codigo}
parCritico :: (Termino -> Termino) -> (Termino, Termino)
           -> (Termino, Termino) -> [(Termino, Termino)]
parCritico c (l1, r1) (l2, r2)
    | sigma1 == Left UNIFICACION = []
    | otherwise = [(sigma r1, sigma (c r2))]
    where sigma1 = unificacion l1 l2
          sigma = aplicaTerm (elimR(sigma1))
          elimR (Right a) = a
\end{codigo}

\item \texttt{(paresCriticos rlista (l,r))} \index{\texttt{paresCriticos}}calcula la lista de todos
  los pares críticos formados al unificar \texttt{rlista} hacia la
  izquierda con un subtérmino \texttt{l}, donde \texttt{l
    $\rightarrow$ r} es una regla. Por ejemplo,

\begin{sesion}
>>> paresCriticos [(V ("x",1), T "f" [V("y",1)])]
                  (V ("x",2), T "g" [V("y",2)])
[]
>>> paresCriticos [(T "f" [V ("x",1)], T "f" [V ("y",1)]),
                   (T "f"[V ("x",1)], T "g" [V ("x",2)]), 
                   (T "g" [V ("x",2)],  T "f" [V ("y",1)])] 
                   (V ("x",1), T "f" [V("y",1)])
[]
>>> paresCriticos [(T "f" [T "g" [V("x",1),V("y",1)],V("z",1)],
                    T "g" [V("x",1),V("z",1)]),
                   (T "g" [V("x",1),V("y",1)],V("x",1))]
                   (T "g" [V("x",1),V("y",1)],
                    T "g" [V("x",1),V("z",1)])
[(T "g" [V ("x",1),V ("z",3)],V ("x",1))]
\end{sesion}

Su código es,

\begin{codigo}
paresCriticos :: [(Termino, Termino)] -> (Termino, Termino)
                 -> [(Termino, Termino)]
paresCriticos rlista (l,r) =
    cps (\x -> x) (renombraTermino m l) (renombraTermino m r)
    where cps _ (V _) _ = []
          cps c ter@(T f ts) r1 = 
                concat(map (parCritico c (ter,r1)) rlista) ++
                 (pcint c f [] ts r1)(pcint c f [] ts r1)
          pcint _ _ _ [] _ = []
          pcint c f ts0 (t:ts1) r2 =
              (cps (\s -> c (T f (ts0++[s]++ts1))) t r2)
                          ++ (pcint c f (ts0++[t]) ts1 r2)
          m = maximum(map
               (\(l3,r3) ->
                maximum(indiceMaximo l3, indiceMaximo r3))
                 rlista) + 1
\end{codigo}

\item \texttt{(paresLista r)} \index{\texttt{paresLista}}es la lista de todos los pares críticos
  de \texttt{r}, incluyendo los triviales. Por ejemplo,

\begin{sesion}
>>> paresLista [(V ("x",1), T "f" [V("y",1)])]
[]
>>> paresLista [(T "f" [T "g" [V("x",1),V("y",1)],V("z",1)],
                 T "g"[V("x",1),V("z",1)]),
                (T "g" [V("x",1),V("y",1)],V("x",1))]
[(T "g" [V ("x",1),V ("z",1)], T "g" [V ("x",1),V ("z",1)]),
 (T "g" [V ("x",1),V ("z",3)], T "f" [V ("x",1),V ("z",3)]),
 (V ("x",1),V ("x",1))]
\end{sesion}

Su código es,

\begin{codigo}
paresLista :: [(Termino, Termino)] -> [(Termino, Termino)]
paresLista r1 = paresLista2 r1 r1
\end{codigo}

\item \texttt{(paresLista2 r1 r2)} \index{\texttt{paresLista2}}es la lista de todos los pares
  críticos al unificar los términos a la izquierda de \texttt{r1} con un
  subtérmino a la izquierda con \texttt{r2}. Por ejemplo,

\begin{sesion}
>>> paresLista2 [(V ("x",1), T "f" [V("y",1)])]
                [(V ("x",1), T "f"[V("y",1)])]
[]
>>> paresLista2 [(T "f" [T "g" [V("x",1),V("y",1)],V("z",1)],
                  T "g"[V("x",1),V("z",1)]),
                 (T "g" [V("x",1),V("y",1)],V("x",1))]
                [(V ("x",1), T "f" [V("y",1)])] 
[]
>>> paresLista2 [(V ("x",1), T "f" [V("y",1)])] 
                [(T "f" [T "g" [V("x",1),V("y",1)],V("z",1)],
                  T "g"[V("x",1),V("z",1)]),
                  (T "g" [V("x",1),V("y",1)],V("x",1))]
[(T "g" [V ("x",3),V ("z",3)],T "f" [V ("y",1)]),
 (T "g" [V ("x",3),V ("z",3)],
  T "f" [T "f" [V ("y",1)],V ("z",3)]),
 (V ("x",3),T "f" [V ("y",1)])]
\end{sesion}

Su código es,

\begin{codigo}
paresLista2 :: [(Termino, Termino)] -> [(Termino, Termino)]
               -> [(Termino, Termino)]
paresLista2 r1 r2 = concat(map(paresCriticos r1) r2)
\end{codigo}

\end{itemize}

\section{Ortogonalidad}

En la sección anterior hemos analizado una técnica para determinar la
confluencia para sistemas terminantes. En esta nos dedicaremos a
analizar el problema para sistemas no terminantes. Para ello la
estudiaremos desde otra perspectiva usando nuevamente los pares
críticos.

\begin{defi}
  Una regla de reescritura $l \rightarrow r$ es lineal por la
  izquierda (respectivamente lineal por la derecha) si ninguna
  variable ocurre dos veces en $l$ (respectivamente $r$). Si una regla
  es lineal por la izquierda y por la derecha diremos que es lineal. 
\end{defi}

Y ampliamos la definición para sistemas de reescritura.

\begin{defi}
  Un sistema de reescritura es lineal por la izquierda
  (respectivamente por la derecha) si todas sus reglas son lineales por
  la izquierda (respectivamente por la derecha)
\end{defi}

Por último, antes de dar el principal resultado de la sección,
definimos una restricción que nos servirá para demostrar la
confluencia fuerte.

\begin{defi}
  Los términos $s_1$ y $s_2$ es fuertemente unible respecto
  $\rightarrow$ si existen términos $t_1$ y $t_2$ tal que
  $s_1 \xrightarrow{=} t_1 \xleftarrow{*} s_2$ y $s_1 \xrightarrow{*}
  t_2 \xleftarrow{=} s_2$
\end{defi}

% AÑADIR DIBUJO DEL ESTILO:
%    /s1\ 
%  t1    t2
%    \s2/

El siguiente resultado se conoce como el lema de confluencia fuerte.

\begin{defi}
  Si $R$ es lineal y cada par crítico de $R$ es fuertemente unible,
  entonces $R$ es fuertemente confluente.
\end{defi}

\begin{demo}
  La prueba se hará reduciendo cada uno de los casos que estudiamos en
  la sección anterior.
  
  \begin{itemize}
    
  \item Caso 1. Es trivial.

  \item Caso 2.1. Se simplifica mediante la restricción de linealidad
    % DIBUJO

  \item Caso 2.2. Como todos los pares críticos son fuertemente
    unible, la instancia de un par crítico es fuertemente unible

  \end{itemize}

\end{demo}

Si un sistema es fuertemente confluente, en particular, también es
confluente. Por tanto acabamos de probar una condición suficiente para la
confluencia.

% EJEMPLO

\begin{defi}
  Un sistema de reescritura es ortogonal si es lineal por la izquierda
  y no tiene pares críticos
\end{defi}

Probar la confluencia para un sistema ortogonal posee el siguiente
problema; al no poseer linealidad por la derecha, en el caso 2.1 es
posible que se aplique mas de una vez la reducción $l_2 \rightarrow r_2$,
y por tanto, lo único que podemos deducir es la confluencia
local. Para evitar este problema, usaremos una relación de reescritura
que posea la propiedad del diamante, que recordemos que el
cumplimiento de esta propiedad implica la confluencia.

Volviendo al caso 2.1, para conseguir la propiedad del diamante
tenemos que arreglar el problema de aplicar repetidas veces la
reducción $l_2 \rightarrow r_2$. Para ello vamos a considerar
reducciones paralelas para los subárboles. Es decir, si tenemos la
figura \ref{fig:pararb1}, consideramos la reducción
$\rightrightarrows$ con la que obtenemos la situación \ref{fig:pararb2}.

\begin{figure}
\centering
\usetikzlibrary{decorations.pathreplacing}
\begin{tikzpicture}

\draw (0,0) node[anchor=north]{}
  -- (10,0) node[anchor=north]{}
  -- (5,5) node[below]{}
  -- cycle
    
       (0,-2) node[anchor=north]{}
  -- (3.33,-2) node[anchor=north]{}
  -- (1.665,0) node[below]{}
  -- cycle
  
       (3.33,-2) node[anchor=north]{}
  -- (6.66,-2) node[anchor=north]{}
  -- (5,0) node[below]{}
  -- cycle
  
        (6.66,-2) node[anchor=north]{}
  -- (10,-2) node[anchor=north]{}
  -- (8.33,0) node[below]{}
  -- cycle
  ;
  

\node[below] at (1.665,-0.8) {\LARGE $\sigma_1 l_1$};
\node[below] at (5,-0.8) {\LARGE $\sigma_2 l_2$};
\node[below] at (8.33,-0.8) {\LARGE $\sigma_3 l_3$};

\end{tikzpicture}
\caption{}
\label{fig:pararb1}
\end{figure}

\begin{figure}
\centering
\usetikzlibrary{decorations.pathreplacing}
\begin{tikzpicture}

\draw (0,0) node[anchor=north]{}
  -- (10,0) node[anchor=north]{}
  -- (5,5) node[below]{}
  -- cycle
    
       (0,-2) node[anchor=north]{}
  -- (3.33,-2) node[anchor=north]{}
  -- (1.665,0) node[below]{}
  -- cycle
  
       (3.33,-2) node[anchor=north]{}
  -- (6.66,-2) node[anchor=north]{}
  -- (5,0) node[below]{}
  -- cycle
  
        (6.66,-2) node[anchor=north]{}
  -- (10,-2) node[anchor=north]{}
  -- (8.33,0) node[below]{}
  -- cycle
  ;

\node(a) at (0,-2) {};
\node (b) at (3.33,-2) {};
\node (c) at (1.665,0) {};
\node(a2) at (3.33,-2) {};
\node (b2) at (6.66,-2) {};
\node (c2) at (5,0) {};
\node(a3) at (6.66,-2) {};
\node (b3) at (10,-2) {};
\node (c3) at (8.33,0) {};
\fill (a.center) -- (b.center) -- (c.center);
\fill (a2.center) -- (b2.center) -- (c2.center);
\fill (a3.center) -- (b3.center) -- (c3.center);

\end{tikzpicture}

\caption{}
\label{fig:pararb2}
\end{figure}

\begin{defi}
  Diremos que un conjunto de posiciones es paralelo si $p||q$ para
  todo $p, q \in P$, donde $P$ es un subconjunto de todas las posibles
  posiciones paralelas.
\end{defi}

Dado un término $t_p$ para cada $p \in P$, definimos la notación,
\[ s[t_p]_{p\in P} : = s[t_{p_1}]_{p_1} \dots [t_{p_n}]_{p_n} \]

\begin{defi}
  Si para cada $p \in P$ tenemos la regla $l_p \rightarrow r_p in R$ y
  para cada sustitución $\sigma_p$ tal que $s|_p = \sigma_p l_p$
  llamaremos reducción paralela a
  $s \rightrightarrows^P_R s[\sigma_p r_p]_{p \in P}$
\end{defi}

A partir de esta definición, enunciamos el teorema de los sistemas
ortogonales

\begin{teor}
  Si $R$ es ortogonal entonces $\rightrightarrows_R$ tiene la
  propiedad del diamante
\end{teor}

\begin{demo}
  Sea $s \rightrightarrows^{P_i} t_i, i = 0,1$. Para esta prueba vamos
  a considerar una partición de $P_i = A_i \cup B_i \cup C$ donde
\begin{itemize}
\item $A_i := \{ p \in P_i | \not \exists q \in P_{1-i}, q \leq p \}$
\item $B_i := \{ p \in P_i | \exists q \in P_{1-i}, q < p \}$
\item $P_0 \cap P_1$
\end{itemize}

Analizamos que elementos hay en cada partición; $A_i$ contiene todas
las posiciones contraídas que estén por encima de
$s \rightrightarrows t_{1-i}$, $B_i$ contiene las que están por debajo
estrictamente de las contraídas en $s \rightrightarrows t_{1-i}$, y
$C$ un conjunto de posiciones contraídas en ambas reducciones.

Como no tenemos pares críticos, todos los $b_i$ son el resultado de
aplicar una regla de reducción hacia la derecha a $a_{1-i}$.

La reducción $t_{1-i} \rightrightarrows u$ tiene las mismas reglas que
las posiciones $B_i$. Lo mismo ocurre para $s \rightrightarrows
t_i$. Nuestro objetivo es llegar a que $t_{1-i} \rightrightarrows u$
también usa las posiciones de $A_i$, para llegar a formar el diamante.

Una reducción en la posición $a_i \in A_i$, no puede ser destruida por
una reducción de algún $b_{1-i} \in B_{1-i}$ porque $R$ es lineal por
la izquierda. Por tanto la reducción $t_{1-i} \rightrightarrows u$ usa
las mismas reglas en las posiciones $A_i$ en la reducción
$s \rightrightarrows t_i$ (que aplicadas en las posiciones de $B_i$).
\end{demo}

Y en particular, tenemos el siguiente corolario.

\begin{coro}
Todo sistema de reescritura ortogonal es confluente.
\end{coro}

%%% Local Variables:
%%% mode: latex
%%% TeX-master: "SRT_en_Haskell"
%%% End:

\chapter{Completación}

Durante los anteriores capítulos hemos estudiado la decibilidad para
el problema de la terminación y confluencia. Hemos probado que ambos
problemas son indecidibles para el caso general. En este capítulo
vamos a estudiar un procedimiento, para que dado un sistema de
reescritura, modificar sus elementos para lograr que sea terminante y
confluente. Formalmente queremos construir un prodecimiento para el
problema de la palabra para un sistema finito de identidades $E$.

Supongamos que tenemos un conjunto
$E := \{ x + 0 \equiv x, x + s(y) \equiv s(x+y) \}$. El sistema de
reescritura generado por $E$ es
$R := \{ x + 0 \rightarrow x, x + s(y) \rightarrow s(x+y)\}$.  Para
probar la termianción usamos el camino de orden lexicográfico inducido
por $>$. La confluencia se verifica pues $R$ no tiene pares críticos.

Para estudiar este conjunto hemos aplicado condiciones suficientes, es
decir, podemos tener un sistema de identidades al que no podamos
aplicar las propiedades que hemos estudiado, y sin embargo sea
terminante y/o confluente. La idea detrás de los algorítmos es añadir
reglas que no cambien la estructura equacional de $E$.

Por ejemplo, supongamos que tenemos un sistema $E$ que es terminante
pero no confluente. Eso quiere decir que existen pares críticos en el
sistema de reescritura $R$ asociado a $E$. Si añadimos alguna regla
adiccional, podemos conseguir que esos pares críticos sean unibles.

Hay que tener en cuenta que al añadir reglas adiccionales se crearán
nuevos pares críticos. Este proceso continuará hasta que no existan
más pares críticos.

\section{Algoritmo de completación}

Vamos a definir un proceso que empieza por un conjunto finito de
identidades $E$, que intenta encontrar un sistema de reescritura
convergente $R$ que sea equivalente a $E$.

\textbf{Algoritmo} \hrulefill

\textbf{Input:} Un conjunto finito $E$ de $\Sigma$-identidades y un
orden de reducción > en $T(\Sigma, V)$.

\textbf{Output:} Si el procedimiento termina; un sistema de
reescritura $R$ finito y convergente que es equivalente a $E$. En caso
contrario devuelve un \texttt{Error}.

\textbf{Inicialización:}
Si existe una identidad $(s \equiv t) \in E$ tal que
$s \not = t, s \not > t$ y $t \not > s$, entonces termina con
\texttt{Error}.
En otro caso,
$i := 0$
$R_0 := \{l \rightarrow r | (l \equiv r) \in E \cup E^{-1} \wedge l > r \}$.

\texttt{Mientras} $R_i \not = R_{i-1}$

$R_{i+1} := R_i$;

Para todo $\langle s, t \rangle \in$ \texttt{ParesCriticos($R_i$)}:

$>$ Reducir $s,t$ a forma normal $s', t'$

$>$ Si $s' \not = t'$, $s' \not > t'$ y $t' \not > s'$ terminar con \texttt{Error}

$>$ Si $s' > t'$ entonces $R_{i+1} := R_{i+1} \cup \{s' \rightarrow t'\}$

$>$ Si $t' > s'$ entonces $R_{i+1} := R_{i+1} \cup \{t' \rightarrow s'\}$

$i := i+1$

Devuelve $R_i$

\hrulefill

Primero eliminamos todas las identidades $t = t$ y calculamos todos
los pares críticos y vamos recorriendo la lista. Cada término
$\langle s, t \rangle$ se reduce a su forma normal $s'$ y $t'$. Si
estos últimos son iguales, el par crítico es unible y se pasa al
siguiente. En caso contrario, transformamos el par crítico en una
regla de reescritura. Se comprueba la terminación de esa regla con
$>$. Si esta última se verifica, se añade al sistema de
reescritura. El proceso termina cuando el sistema no tenga más pares
críticos unibles, ó produzca un \texttt{Error}.

El algoritmo puede producir varias situaciones; que devuelva un
sistema de reescritura en la interacción donde todos los pares
críticos son unibles, que produzca un \texttt{Error}, o que no pare
pues se generen infinitas nuevas reglas.

Por último, vamos a demostrar que el resultado del algoritmo es una
solución para el problema de la palabra.

\begin{teor}
  Sea $E$ un conjunto finito de identidades, y $>$ un orden de
  reducción.

  \begin{enumerate}
  \item Si el resultado del algoritmo aplicado a $E$ y al orden $>$
    termina y devuelve $R_n$, entonces $R_n$ es un sistema de
    reescritura finito convergente equivalente a $E$. Además $R_n$ es
    un procedimiento para el problema de la parada de $E$.

  \item Si el resultado del algoritmo aplicado a $E$ y al orden $>$ no
    termina, entonces $R_\infty := \cap_{i \geq 0} R_i$ es un sistema
    de reescritura infinito convergente que es equivalente a $E$.
  \end{enumerate}

\end{teor}

\begin{demo}
  (1) $R_n$ es finito pues se construye mediante iteraciones finitas,
  donde cada interacción añade reglas $l \rightarrow r$ que satisfacen
  $l > r$. Por tanto por construcción $R_n$ es terminante. Además es
  confluente pues todos sus pares críticos son unibles.

  Por doble inclusión probaremos $\approx_{R_n} = \approx_E$.  La
  primera inclusión se da pues $\approx_{R_0} = \approx_E$ es cierto y
  $R_0 \subseteq R_1 \subseteq \dots \subseteq R_n$.

  La segunda se demuestra mediante inducción para
  $\approx_{R_i} \subseteq \approx_E$ en $i$.  El caso base es
  trivial. Para el paso de la inducción, usaremos que $s \approx_R t$
  es cierto para cada par crítico $\langle s, t \rangle$ de $R_i$, y
  como la forma normal $s'$, $t'$ de $s$ y $t$, verifica
  $s \approx_{R_i} s'$ y $t \approx_{R_i} t'$. Por tanto, tenemos
  $s' \approx_R t'$ que implica $s' \approx_E t'$ por la hipótesis de
  inducción, y que $ \approx_{R_i} \subseteq \approx_E$ es decidible
  por \ref{teor:3.1}.

  (2) $R_{\infty }$ es un sistema infinito donde en cada interacción
  se añade al menos una regla. La equivalencia de $R_{\infty }$ y $E$,
  y la terminación de $R_{\infty }$ se prueba de manera análoga a (1).

  Para probar la confluencia de $R_{\infty }$ vamos a probar que todos
  los pares críticos son unibles. Sea $\langle s,t \rangle$ un par
  crítico entre las reglas $l \rightarrow r$,
  $g \rightarrow d \in R_{\infty }$. Como
  $R_{\infty } = \cup_{i \geq 0} R_i$ y
  $R_0 \subseteq R_1 \subseteq \dots$, esto significa que existe
  $n \geq 0$ tal que $l \rightarrow r$, $g \rightarrow d \in R_n$, que
  implica que $\langle s,t \rangle \in$
  \texttt{paresCriticos($R_n$)}. Este par es unible en $R_n$ ó es
  unible en $R_{n+1}$, pues se le habrá añadido la correspondiente
  regla. Por tanto se da la convergencia.
\end{demo}

\section{Mejorando la completación}

El anterior algoritmo suele generar un gran número de reglas. Solo hay
que ver que para cada par crítico surgen varias. Además hay que hacer
el calculo de los pares críticos en cada paso, lo que hace que el algoritmo sea muy costoso tanto para el tiempo como para la memoria.

Para mejorar el algoritmo de completación, necesitamos varios
resultados que nos aseguren de que el proceso es correcto. En la
siguiente sección hablaremos del algoritmo de Huet. Este algoritmo es
algo mas complejo que el anterior y necesitaremos estos resultados
para comprobar que efectivamente, el proceso funciona.

La idea se basa en extender las reglas de simplificación por un
conjunto de reglas que nos resulten mas beneficiosas. Las reglas son
las siguientes.

\[
  \begin{array}{lcl}
    $\texttt{Deducir} $ & \dfrac{E,R}{E \cup \{s \approx t \}, R} & $Si $ s \leftarrow_R u \rightarrow_R t \\ \\
    $\texttt{Orientar}$ & \dfrac{E \cup \{s \approx' t\},R}{E,R \cup \{s \rightarrow t\}} & $Si $ s > t \\ \\
    $\texttt{Eliminar}$ & \dfrac{E \cup \{s \approx s\}, R}{E,R} & \\ \\
    $\texttt{Simplificar identidad}$   & \dfrac{E \cup \{ s\approx' t\},R}{E \cup \{ u \approx t\},R} & $Si $ s \rightarrow_R\\ \\
    $\texttt{D-simplificar regla}$    & \dfrac{E, R \cup \{s \rightarrow t}{E,R \cup \{s \rightarrow u\}} & $Si $ t \rightarrow_R u   \\ \\
    $\texttt{I-simplificar regla}$    & \dfrac{E, R \cup \{s \rightarrow t}{E \cup \{u \approx t \}, R} & $Si $ s \xrightarrow{\sqsupset}_R u
  \end{array} 
\]

La notación $s \approx' t$ indica que es un par no ordenado, es decir
$s \approx t \in E$ ó $t \approx s \in R$. Denotaremos
$(E,R) \vdash_C (E',R')$ si $(E,R)$ se puede transformar en $(E', R')$
aplicando alguna regla anterior.

Las reglas generan un sistema terminante de reescritura terminante,
pues todas están orientadas con el orden de reducción $>$, de lo que
deducimos el siguiente resultado.

\begin{lema}
 Si $R \subseteq >$ y $(E,R) \vdash_C (E',R')$ entonces $R' \subseteq >$.
\end{lema}

Una consecuencia directa de este lema es dado un sistema de
reescritura $R$ del par $(E,R)$ es terminante si este par se ha
obtenido por uno inicial de la forma $(E_0, \emptyset)$ aplicando las reglas.

El siguiente resultado demuestra que las reglas no cambian la
estructura ecuacional.

\begin{lema}
  $(E_1,R_1) \vdash_C (E_2,R_2)$ implica
  $\approx_{E_1 \cup R_1} = \approx_{E_2 \cup R_2}$.
\end{lema}

\begin{demo}
  La demostración es trivial para las tres primeras reglas, por lo que
  nos centraremos en las tres últimas.
  \begin{itemize}
  \item Para \texttt{Simplificar identidad}, partiendo de que
    $E_1 = E \cup \{s \approx' t\}$, $E_2 = E \cup \{ u \approx t \}$,
    $R_1 = R = R_2$, y $s \rightarrow_R u$, que implica
    $u \approx_{E_1 \cup R_1} t$ y
    $\approx_{E_2 \cup R_2} \subseteq \approx_{E_1 \cup R_1}$. Para la
    otra inclusión partimos de $u \equiv t \in E_2$,
    $s \rightarrow_R u$ y $R = R_2$ implica
    $s \approx_{E_2 \cup R_2} t$ y por tanto,
    $\approx_{E_1 \cup R_1} \subseteq \approx_{E_2 \cup R_2}$.
  \item Para \texttt{D-simplificar regla}, tenemos $E_1 = E = E_2$,
    $R_1 = R \cup \{s \rightarrow t \}, R_2 = R \cup \{s \rightarrow
    u\}$ y $t \rightarrow_R u$. Con $s \rightarrow t \in R_1$,
    $t \rightarrow_R u$ y $R \subseteq R_1$ conseguimos
    $s \approx_{E_1 \cup R_1} u$, y con $s \rightarrow u \in R_2$,
    $t \rightarrow_R u$, y $R \subseteq R_2$ implica
    $s \approx_{E_2 \cup R_2} t$. Esto demuestra
    $\approx_{E_1 \cup R_1} = \approx_{E_2 \cup R_2}$.
  \item Para \texttt{I-simplificar regla} se demuestra de manera análoga
\end{itemize}
\end{demo}

Usando la idea del procedimiento de completación que hemos comentado
antes, formalizamos su definición.

\begin{defi}
  Un procedimiento de completación es un programa que acepta un
  conjunto finito de identidades $E_0$, un orden de reducción $>$, y
  usa las reglas para generar una secuencia (finita o infinita),

  \[
    (E_0,R_0) \vdash_C (E_1,R_1) \vdash_C (E_2,R_2) \vdash_C (E_3,
    R_3) \vdash_C \dots
  \]
  donde $R_0 := \emptyset$.  Además definimos $E_\omega$ como el
  conjunto de identidades persistentes y $R_\omega$ como el conjunto
  de reglas persistentes,

  \[
    E_\omega := \bigcup_{i \geq 0} \bigcap_{j \geq i} E_j
  \]

  \[
R_\omega :=  \bigcup_{i \geq 0} \bigcap_{j \geq i} R_j
  \]
  
\end{defi}

Diremos que la secuencia es un éxito syss $E_\omega = \emptyset$ y
$R_\omega$ es convergente y equivalente a $E_0$. Diremos que falla
syss $E_\omega \not = \emptyset$. Un procedimiento de completación es
correcto syss cada secuencia que no falle, sea un éxito.

Por último, la definición siguiente recoge la importancia de los pares
críticos a la hora de analizar los procedimientos de completación.

\begin{defi}
  Una secuencia de completación se dirá justa syss

  \[
\texttt{paresCriticos}(R_\omega) \subseteq \bigcup_{i\geq 0} E_i
\]
Un procedimiento de completación es justo syss cada secuencia de éxito
es justa.
\end{defi}

Y enunciamos el resultado que nos ayudará a probar que el algoritmo de
Huet resuelve el problema de completación.

\begin{teor}
  Cada procedimiento de completación justo es correcto.
\end{teor}


\section{Algoritmo de Huet}

El algoritmo que vamos a enunciar es una extensión no trivial del
algorítmo básico de completación.

\textbf{Algoritmo de Huet} \hrulefill

\textbf{Input:} Un conjunto finito $E$ de $\Sigma$-identidades y un
orden de reducción > en $T(\Sigma, V)$.

\textbf{Output:} Si el procedimiento termina; un sistema de
reescritura $R$ infinito limitado y convergente que es equivalente a
$E$. En caso contrario devuelve un \texttt{Error}.

\textbf{Inicialización:}
$R_0 := \emptyset, E_O := E, i:= 0;$

\texttt{Mientras} $E_i \not = \emptyset$ ó exista una regla no marcada
en $R_i$

\hspace{0.5cm} \texttt{Mientras} $E_i \not = \emptyset$

\hspace{0.5cm} (a) Escoge una identidad $s \approx t \in E_i$

\hspace{0.5cm} (b) Reduce $s,t$ a una forma $R_i$-normal $s',t'$

\hspace{0.5cm} (c) Si $s' = t'$ entonces,

\hspace{1.4cm} $R_{i+1} := R_i$

\hspace{1.4cm} $E_{i+1} := E_i - \{s \approx t\}$

\hspace{1.4cm} $i := i+1$

\hspace{0.5cm} (d) Si no, si $s' \not > t'$ y $t' \not s'$, entonces
terminar con un output \texttt{Error}.

\hspace{0.5cm} (e) Si no, sea $l,r$ tal que $\{l,r \} = \{ s',t' \}$ y
$l >r$;

\hspace{1.4cm}
$R_{i+1} := \{g \rightarrow d'$ $|$ $g \rightarrow d \in R_i,$ $g$ no
puede reducirse con $l \rightarrow r$, y $d'$ es una forma normal de
$d$ respecto $R_i \cup \{l \rightarrow r\} \}$

\hspace{1.4cm}
$E_{i+1} := (E_i - \{s \approx t \}) \cup \{g' \approx d$ $|$
$g \rightarrow d \in R_i$ y $g$ puede ser reducido a $g'$ con $l \rightarrow r \}$.

\hspace{1.4cm} i := i + 1

\hspace{0.5cm} Si existe una regla no marcada en $R_i$, sea
$l \rightarrow r$ esa regla.

\hspace{1.4cm} $R_{i+1} := R_i$

\hspace{1.4cm} $E_{i+1} := \{s \approx t$ $|$ $\langle s, t \rangle$
es un par crítico de $l \rightarrow r$ ó una regla marcada en $R_i \}$

\hspace{1.4cm} $i := i+1$

\hspace{1.4cm} Marca $l \rightarrow r$.

Devuelve $R_i$

\hrulefill

Comprobamos que el algoritmo efectivamente realiza un procedimiento de
completación sobre $E$

\begin{lema}
El algoritmo de completación de Huet es un procedimiento de completación.
\end{lema}

\begin{demo}
  El objetivo de la prueba es comprobar que cada paso del algoritmo es
  una de las reglas de la anterior sección.
  
  La computación de los pares críticos en el primer \texttt{Mientras}
  se puede realizar mediante \texttt{Deducir}, el paso (b) con
  \texttt{Simplificar identidad}, (c) con \texttt{Eliminar} y (e) con
  \texttt{Orienta}, \texttt{R-simplificar regla} y
  \texttt{L-simplificar regla}.
\end{demo}

Debemos comprobar que si el algoritmo devuelve un \texttt{Error}, el
procedimiento de completación falla

\begin{lema}
Una cadena del algoritmo de Huet falla syss termina con output \texttt{Fail}.
\end{lema}

\begin{demo}
  Si el procedimiento de Huet termina con \texttt{Error}, entonces hay
  una identidad que no puede transformarse por la regla de
  simplificación y orientación, entonces $E_\omega \not = \emptyset$.
  Si el procedimiento no termina con \texttt{Error}, entonces no hay
  identidades y por ello, el bucle interior siempre termina. Por ello
  vamos a asociar cada iteracción del bucle como un multiconjunto de
  pares,

  \[
    K_i := \{ ( \{s,t \},1) | s \approx t \in E_i \} \cup \{ ( \{l,r
    \},0) | l \rightarrow r \in R_i \}
  \]
  Donde la primera componente esta ordenada usando el orden de
  multiconjuntos inducido por el orden de reducción $>$, y la segunda
  componente por el producto del orden lexicográfico.

  Como los multiconjuntos $K_i$ están ordenados usando una extensión
  de este orden de pares, este esta bien fundando y $K_i >
  K_{i+1}$. Por tanto el bucle termina y se demuestra el resultado.
\end{demo}

\section{Implementación del algoritmo de Huet}

La implementación del algoritmo de Huet se realizará por partes.

\begin{itemize}

\item \texttt{(minRegla rl n r1 r2)} \index{minRegla} divide un
  conjunto de reglas en dos, donde uno tiene el tamaño de \texttt{r1}. Por
  ejemplo,
  \begin{sesion}
>>> minRegla (V("x",1),V("x",2)) 2 [(V("y",1),V("y",2))] []
((V ("x",1),V ("x",2)),[(V ("y",1),V ("y",2))])
>>> minRegla (V("x",1),V("x",2)) 1 [(V("y",1),V("y",2))] []
((V ("x",1),V ("x",2)),[(V ("y",1),V ("y",2))])
>>> minRegla (V ("x",1),V ("x",2)) 2
             [(V ("y",1),V ("y",2)),
              (T "f" [V("z",1),V("z",2)], V ("x",2))] []
((V ("x",1),V ("x",2)),
 [(T "f" [V ("z",1),V ("z",2)],V ("x",2)),(V ("y",1),V ("y",2))])
>>> minRegla (V ("x",1),V ("x",2)) 4 [(V ("y",1),V ("y",2)),
             (T "f" [V("z",1),V("z",2)], V ("x",2))] []
((V ("y",1),V ("y",2)),
 [(T "f" [V ("z",1),V ("z",2)],V ("x",2)),(V ("x",1),V ("x",2))])
  \end{sesion}

  Su código es,

  \begin{codigo}
minRegla :: (Termino, Termino) -> Int -> [(Termino, Termino)]
     -> [(Termino, Termino)]
     -> ((Termino, Termino), [(Termino, Termino)])
minRegla rl _ [] r2 = (rl,r2)
minRegla rl n ((l,r):r1) r2
    | m < n = minRegla (l,r) m r1 (rl:r2)
    | otherwise = minRegla rl n r1 ((l,r):r2)
    where m = longitudTerm l + longitudTerm r
  \end{codigo}

\item \texttt{(escogeRegla r)} \index{escogeRegla} calcula la regla
  que hay que aplicar durante el algoritmo de Huet. Por ejemplo,

  \begin{sesion}
>>> escogeRegla [(V("y",1),V("y",2))]
((V ("y",1),V ("y",2)),[])
>>> escogeRegla [(V ("y",1),V ("y",2)),
                 (T "f" [V("z",1),V("z",2)], V ("x",2))]
((V ("y",1),V ("y",2)),[(T "f" [V ("z",1),V ("z",2)],V ("x",2))])
  \end{sesion}

  Su código es,

  \begin{codigo}
escogeRegla :: [(Termino, Termino)]
            -> ((Termino, Termino), [(Termino, Termino)])
escogeRegla [] = error("No se puede aplicar ninguna regla")
escogeRegla ((l,r):r1) = minRegla (l,r)
                         (longitudTerm l + longitudTerm r)
                         r1 []
  \end{codigo}

\item \texttt{(anadeRegla (l,r) e s r)} \index{anadeRegla} es el paso
  (e) del algoritmo de Huet.

  \begin{codigo}
anadeRegla :: (Termino, Termino) -> [(Termino, Termino)]
              -> [Ecuacion] -> [Ecuacion]
              -> ([(Termino, Termino)], [(Termino, Termino)],
                 [(Termino, Termino)])
anadeRegla (l,r) e s r1 = (e2, (l,r):s1, r2)
    where (e1,s1) = simpl l r s r1 s e []
          (e2,r2) = simpl l r s r1 r1 e1 []
simpl
  :: Termino
     -> Termino
     -> [Ecuacion]
     -> [Ecuacion]
     -> [(Termino, Termino)]
     -> [(Termino, Termino)]
     -> [(Termino, Termino)]
     -> ([(Termino, Termino)], [(Termino, Termino)])
simpl _ _ _ _ [] a b = (a,b)
simpl l r s r1 ((g,d):u) e1 u1 
    | g1 == g =
        simpl l r s r1 u e1 ((g, d1):u1)
    | otherwise = simpl l r r1 s u ((g1,d):e1) u1
       where g1 = formaNormal [(l,r)] g
             d1 = formaNormal ((l,r):(r1++s)) d
\end{codigo}

\item \texttt{(orienta ord trip)} \index{orienta} es el segundo bucle
  del algoritmo de Huet.

  \begin{codigo}
orienta :: (Termino -> Termino -> Ordering)
           -> ([(Termino, Termino)],  [(Termino, Termino)],
               [(Termino, Termino)])
           -> ([(Termino, Termino)], [(Termino, Termino)])
orienta _ ([],ss,rr) = (ss,rr)
orienta ord ((s,t):ee,ss,rr)
    | s1 == t1 = orienta ord (ee,ss,rr)
    | ord s1 t1 == GT = orienta ord (anadeRegla (s1,t1)
                                              ee ss rr)
    | ord t1 s1 == GT = orienta ord (anadeRegla (t1,s1)
                                              ee ss rr)
    | otherwise = error("Error con la funcion ord")
    where s1 = formaNormal (rr++ss) s
          t1 = formaNormal (rr++ss) t
  \end{codigo}

\item \texttt{(completa ord ee)} \index{completa} es el resultado de aplicar a ee, el algoritmo de completación de Huet. Por ejemplo,

  \begin{sesion}
>>> let ord = ordenCamLex (ordenPorLista ["a","b","c","x","y","z"])
>>> completa ord [(V("x",1),V("x",1))]
[]
>>> completa ord [(V ("y",1),V ("y",2)),(T "f" [V("y",1),V("y",2)], V ("x",2))]
*** Exception: Error con la funcion ord
>>> completa ord [(T "f" [V("x",1),V("x",2)], V ("x",2))]
[(T "f" [V ("x",1),V ("x",2)],V ("x",2))]
\end{sesion}

Su código es,

\begin{codigo}
completa :: (Termino -> Termino -> Ordering)
            -> [(Termino, Termino)] -> [(Termino, Termino)]
completa ord ee = compl ord (ee,[],[])
compl :: (Termino -> Termino -> Ordering)
         -> ([(Termino, Termino)], [(Termino, Termino)],
             [(Termino, Termino)])
         -> [(Termino, Termino)]
compl ord (ee,ss,tt) = case orienta ord (ee,ss,tt) of
                         ([],rr1)  -> rr1
                         (ss1,rr1) -> compl ord
                                      (paCr,ss2,rl:rr1)
      where (rl, ss2) = escogeRegla ss1
            paCr = (paresLista2 [rl] rr1)++
                   (paresLista2 rr1 [rl])++
                   (paresLista2 [rl] [rl])
\end{codigo}

\end{itemize}

           

%%%%%%%%%%%%%%%%%%%%%%%%%%%%%%%%%%%%%%%%%%%%%%%%%%%
%%Apéndices
%%%%%%%%%%%%%%%%%%%%%%%%%%%%%%%%%%%%%%%%%%%%%%%%%%%

\clearpage
\addappheadtotoc
\appendix


%%% Local Variables:
%%% mode: latex
%%% TeX-master: "SRT_en_Haskell"
%%% End:

%\chapter{Funciones auxiliares para Haskell}

Usaremos las siguientes funciones de la libreria Data.List

\begin{itemize}
\item \texttt{(xs ++ ys)} es la concatenación de \texttt{xs} e \texttt{ys}. Por ejemplo,
        
\begin{sesion}
> [2,5] ++ [3,7,6]
[2,5,3,7,6]
\end{sesion}
        
\item \texttt{(any p xs)} se verifica si algún elemento de \texttt{xs} cumple
  la propiedad \texttt{p}. Por ejemplo,

\begin{sesion}
> any even [3,2,5]
True
> any even [3,1,5]
False
\end{sesion}      

\item \texttt{(all p xs)} se verifica si todos los elementos de \texttt{xs}
  cumplen la propiedad \texttt{p}. Por ejemplo,

\begin{sesion}
> all even [2,6,8]
True
> all even [2,5,8]
False
\end{sesion}      

\item \texttt{(null xs)} se verifica si \texttt{xs} es la lista vacía. Por
  ejemplo, 

\begin{sesion}    
> null []
True
> null [3]
False
\end{sesion}      

\item \texttt{(zip xs ys)} es la lista de pares formado por los
  correspondientes elementos de \texttt{xs} e \texttt{ys}.

\begin{sesion}    
> zip [3,5,2] [4,7]
[(3,4),(5,7)]
> zip [3,5] [4,7,2]
[(3,4),(5,7)]
\end{sesion}              
\end{itemize}

%%% Local Variables:
%%% mode: latex
%%% TeX-master: "SRT_en_Haskell"
%%% End:


%%%%%%%%%%%%%%%%%%%%%%%%%%%%%%%%%%%%%%%%%%%%%%%%%%%%%%%%%%%%%%%%%%%%%%%%%%%%%%%
%%  Bibliografía                                                            %%
%%%%%%%%%%%%%%%%%%%%%%%%%%%%%%%%%%%%%%%%%%%%%%%%%%%%%%%%%%%%%%%%%%%%%%%%%%%%%%%
\nocite{*}
\addcontentsline{toc}{chapter}{Bibliografía}
\bibliographystyle{alpha}
\bibliography{SRT_en_Haskell}

%%%%%%%%%%%%%%%%%%%%%%%%%%%%%%%%%%%%%%%%%%%%%%%%%%%%%%%%%%%%%%%%%%%%%%%%%%%%%%%
%%  Índice                                                                  %%
%%%%%%%%%%%%%%%%%%%%%%%%%%%%%%%%%%%%%%%%%%%%%%%%%%%%%%%%%%%%%%%%%%%%%%%%%%%%%%%

\addcontentsline{toc}{chapter}{Índice de definiciones}

\printindex


















%\begin{codigo}
%       agrupa :: Int -> [a] -> [[a]]
%       agrupa n [] = []
%       agrupa n xs = take n xs : agrupa n (drop n xs)
%\end{codigo}
%\chapter*{English Abstract}
%\addcontentsline{toc}{chapter}{English Abstract}
%\markright{English Abstract}
%\begin{otherlanguage}{english}
%    According to the guidelines, every dissertation %should include a short english abstract at the %beginning. In the abstract, you describe in general %terms what is your dissertation about, the main points %you want to make, and any important consequences that %may arise.
%\end{otherlanguage}


%\chapter{Los enunciados}
%
%\section{Teoremas y demostraciones}
%
%
%\begin{theorem}[Euclides]\label{thm:th1}
%    Esto es un Teorema. Se numeran a partir del 1 en cada capítulo. Como son importantes, tienen un cuadrado rojo al principio. Llevan letra cursiva.
%\end{theorem}
%
%\begin{proof}
%    Esto es la demostración. Al final de la demostración se puede ver un cuadrado rojo similar al de los teoremas. Las demostraciones no llevan letra cursiva.
%\end{proof}
%
%
%\begin{definition}\label{def:1}
%    Esto es una definición. Las definiciones son importantes; también llevan un cuadradito rojo.
%\end{definition}
%
%
%\subsection{Otros enunciados}
%
%
%\begin{remark}
%    Esto es una observación, que dice que $e=mc^{2}$. Como las observaciones no son importantes, no llevan cuadrado rojo, y el tipo de letra no es cursiva.
%\end{remark}
%
%
%\begin{proof}
%    Si la demostración acaba en una fórmula, para poner el cuadrado rojo a la altura de la última formula, hay que usar la orden \verb|\qedhere|, como en este caso:
%    \[
%        e=mc^{2}.\qedhere
%    \]
%
%\end{proof}
%
%
%\begin{corollary}\label{cor:1}
%    Esto es un corolario.
%\end{corollary}
%
%\begin{proposition}\label{pro:1}
%    Esto es una proposición.
%\end{proposition}
%
%\begin{lemma}[Gauss]\label{lem:1}
%    Esto es un lema.
%\end{lemma}
%
%
%\backmatter
%
%\bibliographystyle{acm}
% \biliography{miarchivo} % -> lee miarchivo.bib
% he usado esto: http://www.cs.us.es/~jalonso/cursos/i1m/doc/Funciones_basicas.html#(++)


\end{document}


%%% Local Variables:
%%% mode: latex
%%% TeX-master: t
%%% End:
